\chapter{QAOA : l'algorithme d'optimisation approximative quantique}\label{ch:qaoa}
\iterbanner{V0}

QAOA (\emph{Quantum Approximate Optimization Algorithm}, Farhi, Goldstone et Gutmann, 2014) est l'algorithme que la note d'intention prévoit d'exécuter sur simulateur. Il réunit tout ce qui précède : un hamiltonien de coût diagonal (ch.~\ref{ch:qubo}), un mélangeur qui ne commute pas avec lui (ch.~\ref{ch:algebre}), une alternance à la Trotter (ch.~\ref{ch:hamiltoniens}) et un principe variationnel réglé par un ordinateur classique.

\section{De l'adiabatique à QAOA}
Discrétisons le recuit adiabatique $H(s)=(1-s)H_B+sH_C$ en $p$ pas, et séparons chaque pas par Trotter :
\[e^{-iH(s_k)\Delta t}\approx e^{-i(1-s_k)\Delta t\,H_B}\ e^{-is_k\Delta t\,H_C}.\]
On obtient une alternance de deux types d'opérations, dont les « durées » suivent un programme fixé : $s_k\Delta t$ croissant pour le coût, $(1-s_k)\Delta t$ décroissant pour le mélangeur. \textbf{L'idée de QAOA : ne pas imposer ces durées, mais les optimiser.}

\section{Définition}
\begin{definition}[ (QAOA de profondeur $p$)]
Avec $H_C$ l'hamiltonien d'Ising du problème et le mélangeur $H_M=\sum_kX_k$, on prépare
\[\ket{\boldsymbol\gamma,\boldsymbol\beta}=e^{-i\beta_pH_M}e^{-i\gamma_pH_C}\cdots e^{-i\beta_1H_M}e^{-i\gamma_1H_C}\ \ket+^{\otimes n},\]
puis on cherche les $2p$ angles qui minimisent
\[F_p(\boldsymbol\gamma,\boldsymbol\beta)=\bra{\boldsymbol\gamma,\boldsymbol\beta}H_C\ket{\boldsymbol\gamma,\boldsymbol\beta}.\]
On mesure ensuite l'état optimal en base de calcul : chaque shot donne une chaîne de bits candidate.
\end{definition}
\paragraph{Ce que fait chaque étage, intuitivement.}
\begin{itemize}
\item $\ket+^{\otimes n}$ : toutes les chaînes ont la même amplitude, $2^{-n/2}$ ;
\item $e^{-i\gamma H_C}$ : chaque chaîne $x$ reçoit une phase $e^{-i\gamma E(x)}$. Le coût est « écrit » dans les phases, les probabilités sont inchangées ;
\item $e^{-i\beta H_M}$ : chaque qubit tourne autour de $x$. Les amplitudes de chaînes voisines (un bit de différence) se mélangent et \textbf{interfèrent}. Selon leurs phases, elles se renforcent ou s'annulent. Bien réglé, ce mélange concentre la probabilité sur les chaînes de basse énergie.
\end{itemize}
\paragraph{Garanties.} $F_p$ décroît (au sens large) avec $p$, et quand $p\to\infty$, QAOA contient le recuit adiabatique : il atteint l'optimum. À $p$ fini, c'est une heuristique.

\section{Le circuit}
On utilise $e^{-i\gamma h_iZ_i}=R_Z(2\gamma h_i)$, $e^{-i\gamma J_{ij}Z_iZ_j}=R_{ZZ}(2\gamma J_{ij})$ et $e^{-i\beta X_k}=R_X(2\beta)$. Tous les termes de $H_C$ commutent entre eux (ils sont tous diagonaux) : $e^{-i\gamma H_C}$ se factorise donc \emph{exactement}, sans erreur de Trotter.

\begin{figure}[htbp]
\centering
\begin{tikzpicture}[x=1cm,y=0.95cm]
\foreach \q in {0,1,2}{\draw[cink] (0,-\q) -- (11.6,-\q); \node[anchor=east,font=\small] at (0,-\q) {$q_\q:\ \ket0$};}
\foreach \q in {0,1,2}{\node[gate] at (0.6,-\q) {$H$};}
\draw[cmuted,dashed,rounded corners] (1.15,0.55) rectangle (7.55,-2.55);
\node[font=\scriptsize,cmuted] at (4.35,0.78) {coût $e^{-i\gamma_kH_C}$};
\node[gate,font=\footnotesize,minimum width=17mm] at (2.25,-1) {$R_Z(2\gamma_kh_1)$};
\node[draw=cink,fill=white,minimum width=20mm,minimum height=14mm,font=\footnotesize] at (4.3,-0.5) {$R_{ZZ}(2\gamma_kJ_{01})$};
\node[draw=cink,fill=white,minimum width=20mm,minimum height=14mm,font=\footnotesize] at (6.45,-1.5) {$R_{ZZ}(2\gamma_kJ_{12})$};
\draw[cmuted,dashed,rounded corners] (7.75,0.55) rectangle (9.65,-2.55);
\node[font=\scriptsize,cmuted] at (8.7,0.78) {mélangeur};
\foreach \q in {0,1,2}{\node[gate,font=\footnotesize,minimum width=14mm] at (8.7,-\q) {$R_X(2\beta_k)$};}
\draw[decorate,decoration={brace,amplitude=5pt,mirror}] (1.15,-2.75) -- (9.65,-2.75) node[midway,below=6pt,font=\small] {répété pour $k=1,\dots,p$};
\foreach \q in {0,1,2}{\meter{10.7}{-\q}}
\end{tikzpicture}
\caption{Circuit QAOA pour le QUBO jouet à 3 qubits ($h_0=h_2=0$, $h_1=-\frac12$, $J_{01}=J_{12}=\frac12$). Chaque $R_{ZZ}$ se décompose en 2 CNOT (figure~\ref{fig:rzz}).}\label{fig:circuit-qaoa}
\end{figure}

\paragraph{Ressources par couche.} Pour $n$ qubits et $|E|$ couplages non nuls : $n$ portes $R_Z$, $|E|$ portes $R_{ZZ}$ (donc $2|E|$ CNOT) et $n$ portes $R_X$. Au total, $p$ couches, $n$ portes $H$ initiales et $n$ mesures. Le nombre de qubits vaut $n$, le nombre de bits du QUBO.

\section{La boucle hybride quantique--classique}
\begin{figure}[htbp]
\centering
\begin{tikzpicture}[node distance=10mm and 16mm]
\node[flow,fill=cblue!10,minimum width=4.2cm,minimum height=1.5cm] (q) {\textbf{Processeur quantique}\\ (ou simulateur)\\ prépare $\ket{\boldsymbol\gamma,\boldsymbol\beta}$, mesure $N$ shots};
\node[flow,right=of q,minimum width=4.2cm,minimum height=1.5cm] (c) {\textbf{Optimiseur classique}\\ (COBYLA, Nelder-Mead)\\ propose de nouveaux $(\boldsymbol\gamma,\boldsymbol\beta)$};
\draw[arr] (q.north east) to[bend left=20] node[above,font=\small] {estimation de $F_p$} (c.north west);
\draw[arr] (c.south west) to[bend left=20] node[below,font=\small] {$(\boldsymbol\gamma,\boldsymbol\beta)$} (q.south east);
\node[flow,below=14mm of q,minimum width=4.2cm] (s) {échantillonnage final\\ $\to$ meilleure chaîne $x^*$};
\draw[arr] (q) -- node[left,font=\small] {à convergence} (s);
\end{tikzpicture}
\caption{QAOA est un algorithme \emph{variationnel} : le circuit quantique évalue la fonction, l'ordinateur classique l'optimise. Les $2p$ angles sont les seuls paramètres.}\label{fig:boucle-hybride}
\end{figure}

\section{Exemple entièrement calculé : un qubit}
Prenons $H_C=Z$ : les énergies sont $+1$ pour $\ket0$ et $-1$ pour $\ket1$, donc l'optimum est $\ket1$. Avec $p=1$, suivons le point sur la sphère de Bloch :
\begin{enumerate}
\item $\ket+$ est le point $(1,0,0)$ ;
\item $e^{-i\gamma Z}=R_Z(2\gamma)$ le fait tourner de $2\gamma$ autour de $z$ : $(\cos2\gamma,\ \sin2\gamma,\ 0)$ ;
\item $e^{-i\beta X}=R_X(2\beta)$ le fait tourner de $2\beta$ autour de $x$, ce qui donne une composante $z=\sin2\gamma\,\sin2\beta$.
\end{enumerate}
Comme $\expval Z$ est la coordonnée $z$ :
\[F_1(\gamma,\beta)=\sin2\gamma\ \sin2\beta,\qquad \min F_1=-1\ \text{ en }(\gamma,\beta)=(\tfrac\pi4,-\tfrac\pi4)\ \Rightarrow\ \text{on mesure }\ket1\text{ avec certitude}.\]
Le coût a d'abord fait tourner l'état \emph{dans l'équateur} (une phase), puis le mélangeur a converti cette phase en déplacement vers le pôle de basse énergie. C'est le mécanisme général, sur $2^n$ amplitudes au lieu d'une.

\section{Résultats sur le QUBO jouet (3 qubits)}
\begin{figure}[htbp]
\centering
% figure générée par tools/make_figures.py : paysage_qaoa_p1
\begin{tikzpicture}
\begin{axis}[width=0.78\linewidth,height=6cm,view={0}{90},xmin=0,xmax=6.2832,ymin=0,ymax=3.1416,
  xlabel={$\gamma$},ylabel={$\beta$},xtick={0,1.5708,3.1416,4.7124,6.2832},
  xticklabels={$0$,$\pi/2$,$\pi$,$3\pi/2$,$2\pi$},ytick={0,1.5708,3.1416},yticklabels={$0$,$\pi/2$,$\pi$},
  colormap={seq}{color=(cblue!85!black) color=(cblue!45) color=(white)},
  colorbar,colorbar style={ylabel={$F_1(\gamma,\beta)$}},axis on top,enlargelimits=false]
\addplot3[surf,shader=interp,mesh/rows=37] coordinates {
(0.0000,0.0000,-0.5000) (0.0873,0.0000,-0.5000) (0.1745,0.0000,-0.5000) (0.2618,0.0000,-0.5000) (0.3491,0.0000,-0.5000) (0.4363,0.0000,-0.5000) (0.5236,0.0000,-0.5000) (0.6109,0.0000,-0.5000) (0.6981,0.0000,-0.5000) (0.7854,0.0000,-0.5000) (0.8727,0.0000,-0.5000) (0.9599,0.0000,-0.5000) (1.0472,0.0000,-0.5000) (1.1345,0.0000,-0.5000) (1.2217,0.0000,-0.5000) (1.3090,0.0000,-0.5000) (1.3963,0.0000,-0.5000) (1.4835,0.0000,-0.5000) (1.5708,0.0000,-0.5000) (1.6581,0.0000,-0.5000) (1.7453,0.0000,-0.5000) (1.8326,0.0000,-0.5000) (1.9199,0.0000,-0.5000) (2.0071,0.0000,-0.5000) (2.0944,0.0000,-0.5000) (2.1817,0.0000,-0.5000) (2.2689,0.0000,-0.5000) (2.3562,0.0000,-0.5000) (2.4435,0.0000,-0.5000) (2.5307,0.0000,-0.5000) (2.6180,0.0000,-0.5000) (2.7053,0.0000,-0.5000) (2.7925,0.0000,-0.5000) (2.8798,0.0000,-0.5000) (2.9671,0.0000,-0.5000) (3.0543,0.0000,-0.5000) (3.1416,0.0000,-0.5000) (3.2289,0.0000,-0.5000) (3.3161,0.0000,-0.5000) (3.4034,0.0000,-0.5000) (3.4907,0.0000,-0.5000) (3.5779,0.0000,-0.5000) (3.6652,0.0000,-0.5000) (3.7525,0.0000,-0.5000) (3.8397,0.0000,-0.5000) (3.9270,0.0000,-0.5000) (4.0143,0.0000,-0.5000) (4.1015,0.0000,-0.5000) (4.1888,0.0000,-0.5000) (4.2761,0.0000,-0.5000) (4.3633,0.0000,-0.5000) (4.4506,0.0000,-0.5000) (4.5379,0.0000,-0.5000) (4.6251,0.0000,-0.5000) (4.7124,0.0000,-0.5000) (4.7997,0.0000,-0.5000) (4.8869,0.0000,-0.5000) (4.9742,0.0000,-0.5000) (5.0615,0.0000,-0.5000) (5.1487,0.0000,-0.5000) (5.2360,0.0000,-0.5000) (5.3233,0.0000,-0.5000) (5.4105,0.0000,-0.5000) (5.4978,0.0000,-0.5000) (5.5851,0.0000,-0.5000) (5.6723,0.0000,-0.5000) (5.7596,0.0000,-0.5000) (5.8469,0.0000,-0.5000) (5.9341,0.0000,-0.5000) (6.0214,0.0000,-0.5000) (6.1087,0.0000,-0.5000) (6.1959,0.0000,-0.5000) (6.2832,0.0000,-0.5000)
(0.0000,0.0873,-0.5000) (0.0873,0.0873,-0.4628) (0.1745,0.0873,-0.4269) (0.2618,0.0873,-0.3935) (0.3491,0.0873,-0.3636) (0.4363,0.0873,-0.3382) (0.5236,0.0873,-0.3178) (0.6109,0.0873,-0.3027) (0.6981,0.0873,-0.2928) (0.7854,0.0873,-0.2879) (0.8727,0.0873,-0.2874) (0.9599,0.0873,-0.2904) (1.0472,0.0873,-0.2961) (1.1345,0.0873,-0.3033) (1.2217,0.0873,-0.3110) (1.3090,0.0873,-0.3181) (1.3963,0.0873,-0.3239) (1.4835,0.0873,-0.3277) (1.5708,0.0873,-0.3290) (1.6581,0.0873,-0.3277) (1.7453,0.0873,-0.3239) (1.8326,0.0873,-0.3181) (1.9199,0.0873,-0.3110) (2.0071,0.0873,-0.3033) (2.0944,0.0873,-0.2961) (2.1817,0.0873,-0.2904) (2.2689,0.0873,-0.2874) (2.3562,0.0873,-0.2879) (2.4435,0.0873,-0.2928) (2.5307,0.0873,-0.3027) (2.6180,0.0873,-0.3178) (2.7053,0.0873,-0.3382) (2.7925,0.0873,-0.3636) (2.8798,0.0873,-0.3935) (2.9671,0.0873,-0.4269) (3.0543,0.0873,-0.4628) (3.1416,0.0873,-0.5000) (3.2289,0.0873,-0.5372) (3.3161,0.0873,-0.5731) (3.4034,0.0873,-0.6065) (3.4907,0.0873,-0.6364) (3.5779,0.0873,-0.6618) (3.6652,0.0873,-0.6822) (3.7525,0.0873,-0.6973) (3.8397,0.0873,-0.7072) (3.9270,0.0873,-0.7121) (4.0143,0.0873,-0.7126) (4.1015,0.0873,-0.7096) (4.1888,0.0873,-0.7039) (4.2761,0.0873,-0.6967) (4.3633,0.0873,-0.6890) (4.4506,0.0873,-0.6819) (4.5379,0.0873,-0.6761) (4.6251,0.0873,-0.6723) (4.7124,0.0873,-0.6710) (4.7997,0.0873,-0.6723) (4.8869,0.0873,-0.6761) (4.9742,0.0873,-0.6819) (5.0615,0.0873,-0.6890) (5.1487,0.0873,-0.6967) (5.2360,0.0873,-0.7039) (5.3233,0.0873,-0.7096) (5.4105,0.0873,-0.7126) (5.4978,0.0873,-0.7121) (5.5851,0.0873,-0.7072) (5.6723,0.0873,-0.6973) (5.7596,0.0873,-0.6822) (5.8469,0.0873,-0.6618) (5.9341,0.0873,-0.6364) (6.0214,0.0873,-0.6065) (6.1087,0.0873,-0.5731) (6.1959,0.0873,-0.5372) (6.2832,0.0873,-0.5000)
(0.0000,0.1745,-0.5000) (0.0873,0.1745,-0.4294) (0.1745,0.1745,-0.3613) (0.2618,0.1745,-0.2979) (0.3491,0.1745,-0.2414) (0.4363,0.1745,-0.1932) (0.5236,0.1745,-0.1547) (0.6109,0.1745,-0.1261) (0.6981,0.1745,-0.1077) (0.7854,0.1745,-0.0986) (0.8727,0.1745,-0.0979) (0.9599,0.1745,-0.1040) (1.0472,0.1745,-0.1151) (1.1345,0.1745,-0.1290) (1.2217,0.1745,-0.1439) (1.3090,0.1745,-0.1577) (1.3963,0.1745,-0.1689) (1.4835,0.1745,-0.1761) (1.5708,0.1745,-0.1786) (1.6581,0.1745,-0.1761) (1.7453,0.1745,-0.1689) (1.8326,0.1745,-0.1577) (1.9199,0.1745,-0.1439) (2.0071,0.1745,-0.1290) (2.0944,0.1745,-0.1151) (2.1817,0.1745,-0.1040) (2.2689,0.1745,-0.0979) (2.3562,0.1745,-0.0986) (2.4435,0.1745,-0.1077) (2.5307,0.1745,-0.1261) (2.6180,0.1745,-0.1547) (2.7053,0.1745,-0.1932) (2.7925,0.1745,-0.2414) (2.8798,0.1745,-0.2979) (2.9671,0.1745,-0.3613) (3.0543,0.1745,-0.4294) (3.1416,0.1745,-0.5000) (3.2289,0.1745,-0.5706) (3.3161,0.1745,-0.6387) (3.4034,0.1745,-0.7021) (3.4907,0.1745,-0.7586) (3.5779,0.1745,-0.8068) (3.6652,0.1745,-0.8453) (3.7525,0.1745,-0.8739) (3.8397,0.1745,-0.8923) (3.9270,0.1745,-0.9014) (4.0143,0.1745,-0.9021) (4.1015,0.1745,-0.8960) (4.1888,0.1745,-0.8849) (4.2761,0.1745,-0.8710) (4.3633,0.1745,-0.8561) (4.4506,0.1745,-0.8423) (4.5379,0.1745,-0.8311) (4.6251,0.1745,-0.8239) (4.7124,0.1745,-0.8214) (4.7997,0.1745,-0.8239) (4.8869,0.1745,-0.8311) (4.9742,0.1745,-0.8423) (5.0615,0.1745,-0.8561) (5.1487,0.1745,-0.8710) (5.2360,0.1745,-0.8849) (5.3233,0.1745,-0.8960) (5.4105,0.1745,-0.9021) (5.4978,0.1745,-0.9014) (5.5851,0.1745,-0.8923) (5.6723,0.1745,-0.8739) (5.7596,0.1745,-0.8453) (5.8469,0.1745,-0.8068) (5.9341,0.1745,-0.7586) (6.0214,0.1745,-0.7021) (6.1087,0.1745,-0.6387) (6.1959,0.1745,-0.5706) (6.2832,0.1745,-0.5000)
(0.0000,0.2618,-0.5000) (0.0873,0.2618,-0.4032) (0.1745,0.2618,-0.3098) (0.2618,0.2618,-0.2230) (0.3491,0.2618,-0.1456) (0.4363,0.2618,-0.0799) (0.5236,0.2618,-0.0274) (0.6109,0.2618,0.0112) (0.6981,0.2618,0.0360) (0.7854,0.2618,0.0477) (0.8727,0.2618,0.0479) (0.9599,0.2618,0.0388) (1.0472,0.2618,0.0229) (1.1345,0.2618,0.0030) (1.2217,0.2618,-0.0180) (1.3090,0.2618,-0.0375) (1.3963,0.2618,-0.0533) (1.4835,0.2618,-0.0635) (1.5708,0.2618,-0.0670) (1.6581,0.2618,-0.0635) (1.7453,0.2618,-0.0533) (1.8326,0.2618,-0.0375) (1.9199,0.2618,-0.0180) (2.0071,0.2618,0.0030) (2.0944,0.2618,0.0229) (2.1817,0.2618,0.0388) (2.2689,0.2618,0.0479) (2.3562,0.2618,0.0477) (2.4435,0.2618,0.0360) (2.5307,0.2618,0.0112) (2.6180,0.2618,-0.0274) (2.7053,0.2618,-0.0799) (2.7925,0.2618,-0.1456) (2.8798,0.2618,-0.2230) (2.9671,0.2618,-0.3098) (3.0543,0.2618,-0.4032) (3.1416,0.2618,-0.5000) (3.2289,0.2618,-0.5968) (3.3161,0.2618,-0.6902) (3.4034,0.2618,-0.7770) (3.4907,0.2618,-0.8544) (3.5779,0.2618,-0.9201) (3.6652,0.2618,-0.9726) (3.7525,0.2618,-1.0112) (3.8397,0.2618,-1.0360) (3.9270,0.2618,-1.0477) (4.0143,0.2618,-1.0479) (4.1015,0.2618,-1.0388) (4.1888,0.2618,-1.0229) (4.2761,0.2618,-1.0030) (4.3633,0.2618,-0.9820) (4.4506,0.2618,-0.9625) (4.5379,0.2618,-0.9467) (4.6251,0.2618,-0.9365) (4.7124,0.2618,-0.9330) (4.7997,0.2618,-0.9365) (4.8869,0.2618,-0.9467) (4.9742,0.2618,-0.9625) (5.0615,0.2618,-0.9820) (5.1487,0.2618,-1.0030) (5.2360,0.2618,-1.0229) (5.3233,0.2618,-1.0388) (5.4105,0.2618,-1.0479) (5.4978,0.2618,-1.0477) (5.5851,0.2618,-1.0360) (5.6723,0.2618,-1.0112) (5.7596,0.2618,-0.9726) (5.8469,0.2618,-0.9201) (5.9341,0.2618,-0.8544) (6.0214,0.2618,-0.7770) (6.1087,0.2618,-0.6902) (6.1959,0.2618,-0.5968) (6.2832,0.2618,-0.5000)
(0.0000,0.3491,-0.5000) (0.0873,0.3491,-0.3867) (0.1745,0.3491,-0.2774) (0.2618,0.3491,-0.1760) (0.3491,0.3491,-0.0858) (0.4363,0.3491,-0.0094) (0.5236,0.3491,0.0514) (0.6109,0.3491,0.0956) (0.6981,0.3491,0.1235) (0.7854,0.3491,0.1359) (0.8727,0.3491,0.1348) (0.9599,0.3491,0.1227) (1.0472,0.3491,0.1026) (1.1345,0.3491,0.0780) (1.2217,0.3491,0.0522) (1.3090,0.3491,0.0283) (1.3963,0.3491,0.0091) (1.4835,0.3491,-0.0033) (1.5708,0.3491,-0.0076) (1.6581,0.3491,-0.0033) (1.7453,0.3491,0.0091) (1.8326,0.3491,0.0283) (1.9199,0.3491,0.0522) (2.0071,0.3491,0.0780) (2.0944,0.3491,0.1026) (2.1817,0.3491,0.1227) (2.2689,0.3491,0.1348) (2.3562,0.3491,0.1359) (2.4435,0.3491,0.1235) (2.5307,0.3491,0.0956) (2.6180,0.3491,0.0514) (2.7053,0.3491,-0.0094) (2.7925,0.3491,-0.0858) (2.8798,0.3491,-0.1760) (2.9671,0.3491,-0.2774) (3.0543,0.3491,-0.3867) (3.1416,0.3491,-0.5000) (3.2289,0.3491,-0.6133) (3.3161,0.3491,-0.7226) (3.4034,0.3491,-0.8240) (3.4907,0.3491,-0.9142) (3.5779,0.3491,-0.9906) (3.6652,0.3491,-1.0514) (3.7525,0.3491,-1.0956) (3.8397,0.3491,-1.1235) (3.9270,0.3491,-1.1359) (4.0143,0.3491,-1.1348) (4.1015,0.3491,-1.1227) (4.1888,0.3491,-1.1026) (4.2761,0.3491,-1.0780) (4.3633,0.3491,-1.0522) (4.4506,0.3491,-1.0283) (4.5379,0.3491,-1.0091) (4.6251,0.3491,-0.9967) (4.7124,0.3491,-0.9924) (4.7997,0.3491,-0.9967) (4.8869,0.3491,-1.0091) (4.9742,0.3491,-1.0283) (5.0615,0.3491,-1.0522) (5.1487,0.3491,-1.0780) (5.2360,0.3491,-1.1026) (5.3233,0.3491,-1.1227) (5.4105,0.3491,-1.1348) (5.4978,0.3491,-1.1359) (5.5851,0.3491,-1.1235) (5.6723,0.3491,-1.0956) (5.7596,0.3491,-1.0514) (5.8469,0.3491,-0.9906) (5.9341,0.3491,-0.9142) (6.0214,0.3491,-0.8240) (6.1087,0.3491,-0.7226) (6.1959,0.3491,-0.6133) (6.2832,0.3491,-0.5000)
(0.0000,0.4363,-0.5000) (0.0873,0.4363,-0.3814) (0.1745,0.4363,-0.2671) (0.2618,0.4363,-0.1612) (0.3491,0.4363,-0.0672) (0.4363,0.4363,0.0120) (0.5236,0.4363,0.0745) (0.6109,0.4363,0.1194) (0.6981,0.4363,0.1467) (0.7854,0.4363,0.1577) (0.8727,0.4363,0.1543) (0.9599,0.4363,0.1393) (1.0472,0.4363,0.1160) (1.1345,0.4363,0.0880) (1.2217,0.4363,0.0589) (1.3090,0.4363,0.0323) (1.3963,0.4363,0.0109) (1.4835,0.4363,-0.0028) (1.5708,0.4363,-0.0076) (1.6581,0.4363,-0.0028) (1.7453,0.4363,0.0109) (1.8326,0.4363,0.0323) (1.9199,0.4363,0.0589) (2.0071,0.4363,0.0880) (2.0944,0.4363,0.1160) (2.1817,0.4363,0.1393) (2.2689,0.4363,0.1543) (2.3562,0.4363,0.1577) (2.4435,0.4363,0.1467) (2.5307,0.4363,0.1194) (2.6180,0.4363,0.0745) (2.7053,0.4363,0.0120) (2.7925,0.4363,-0.0672) (2.8798,0.4363,-0.1612) (2.9671,0.4363,-0.2671) (3.0543,0.4363,-0.3814) (3.1416,0.4363,-0.5000) (3.2289,0.4363,-0.6186) (3.3161,0.4363,-0.7329) (3.4034,0.4363,-0.8388) (3.4907,0.4363,-0.9328) (3.5779,0.4363,-1.0120) (3.6652,0.4363,-1.0745) (3.7525,0.4363,-1.1194) (3.8397,0.4363,-1.1467) (3.9270,0.4363,-1.1577) (4.0143,0.4363,-1.1543) (4.1015,0.4363,-1.1393) (4.1888,0.4363,-1.1160) (4.2761,0.4363,-1.0880) (4.3633,0.4363,-1.0589) (4.4506,0.4363,-1.0323) (4.5379,0.4363,-1.0109) (4.6251,0.4363,-0.9972) (4.7124,0.4363,-0.9924) (4.7997,0.4363,-0.9972) (4.8869,0.4363,-1.0109) (4.9742,0.4363,-1.0323) (5.0615,0.4363,-1.0589) (5.1487,0.4363,-1.0880) (5.2360,0.4363,-1.1160) (5.3233,0.4363,-1.1393) (5.4105,0.4363,-1.1543) (5.4978,0.4363,-1.1577) (5.5851,0.4363,-1.1467) (5.6723,0.4363,-1.1194) (5.7596,0.4363,-1.0745) (5.8469,0.4363,-1.0120) (5.9341,0.4363,-0.9328) (6.0214,0.4363,-0.8388) (6.1087,0.4363,-0.7329) (6.1959,0.4363,-0.6186) (6.2832,0.4363,-0.5000)
(0.0000,0.5236,-0.5000) (0.0873,0.5236,-0.3874) (0.1745,0.5236,-0.2790) (0.2618,0.5236,-0.1788) (0.3491,0.5236,-0.0904) (0.4363,0.5236,-0.0164) (0.5236,0.5236,0.0413) (0.6109,0.5236,0.0817) (0.6981,0.5236,0.1050) (0.7854,0.5236,0.1124) (0.8727,0.5236,0.1058) (0.9599,0.5236,0.0881) (1.0472,0.5236,0.0625) (1.1345,0.5236,0.0326) (1.2217,0.5236,0.0021) (1.3090,0.5236,-0.0257) (1.3963,0.5236,-0.0478) (1.4835,0.5236,-0.0621) (1.5708,0.5236,-0.0670) (1.6581,0.5236,-0.0621) (1.7453,0.5236,-0.0478) (1.8326,0.5236,-0.0257) (1.9199,0.5236,0.0021) (2.0071,0.5236,0.0326) (2.0944,0.5236,0.0625) (2.1817,0.5236,0.0881) (2.2689,0.5236,0.1058) (2.3562,0.5236,0.1124) (2.4435,0.5236,0.1050) (2.5307,0.5236,0.0817) (2.6180,0.5236,0.0413) (2.7053,0.5236,-0.0164) (2.7925,0.5236,-0.0904) (2.8798,0.5236,-0.1788) (2.9671,0.5236,-0.2790) (3.0543,0.5236,-0.3874) (3.1416,0.5236,-0.5000) (3.2289,0.5236,-0.6126) (3.3161,0.5236,-0.7210) (3.4034,0.5236,-0.8212) (3.4907,0.5236,-0.9096) (3.5779,0.5236,-0.9836) (3.6652,0.5236,-1.0413) (3.7525,0.5236,-1.0817) (3.8397,0.5236,-1.1050) (3.9270,0.5236,-1.1124) (4.0143,0.5236,-1.1058) (4.1015,0.5236,-1.0881) (4.1888,0.5236,-1.0625) (4.2761,0.5236,-1.0326) (4.3633,0.5236,-1.0021) (4.4506,0.5236,-0.9743) (4.5379,0.5236,-0.9522) (4.6251,0.5236,-0.9379) (4.7124,0.5236,-0.9330) (4.7997,0.5236,-0.9379) (4.8869,0.5236,-0.9522) (4.9742,0.5236,-0.9743) (5.0615,0.5236,-1.0021) (5.1487,0.5236,-1.0326) (5.2360,0.5236,-1.0625) (5.3233,0.5236,-1.0881) (5.4105,0.5236,-1.1058) (5.4978,0.5236,-1.1124) (5.5851,0.5236,-1.1050) (5.6723,0.5236,-1.0817) (5.7596,0.5236,-1.0413) (5.8469,0.5236,-0.9836) (5.9341,0.5236,-0.9096) (6.0214,0.5236,-0.8212) (6.1087,0.5236,-0.7210) (6.1959,0.5236,-0.6126) (6.2832,0.5236,-0.5000)
(0.0000,0.6109,-0.5000) (0.0873,0.6109,-0.4036) (0.1745,0.6109,-0.3109) (0.2618,0.6109,-0.2257) (0.3491,0.6109,-0.1511) (0.4363,0.6109,-0.0895) (0.5236,0.6109,-0.0426) (0.6109,0.6109,-0.0111) (0.6981,0.6109,0.0050) (0.7854,0.6109,0.0070) (0.8727,0.6109,-0.0034) (0.9599,0.6109,-0.0235) (1.0472,0.6109,-0.0504) (1.1345,0.6109,-0.0806) (1.2217,0.6109,-0.1110) (1.3090,0.6109,-0.1384) (1.3963,0.6109,-0.1600) (1.4835,0.6109,-0.1738) (1.5708,0.6109,-0.1786) (1.6581,0.6109,-0.1738) (1.7453,0.6109,-0.1600) (1.8326,0.6109,-0.1384) (1.9199,0.6109,-0.1110) (2.0071,0.6109,-0.0806) (2.0944,0.6109,-0.0504) (2.1817,0.6109,-0.0235) (2.2689,0.6109,-0.0034) (2.3562,0.6109,0.0070) (2.4435,0.6109,0.0050) (2.5307,0.6109,-0.0111) (2.6180,0.6109,-0.0426) (2.7053,0.6109,-0.0895) (2.7925,0.6109,-0.1511) (2.8798,0.6109,-0.2257) (2.9671,0.6109,-0.3109) (3.0543,0.6109,-0.4036) (3.1416,0.6109,-0.5000) (3.2289,0.6109,-0.5964) (3.3161,0.6109,-0.6891) (3.4034,0.6109,-0.7743) (3.4907,0.6109,-0.8489) (3.5779,0.6109,-0.9105) (3.6652,0.6109,-0.9574) (3.7525,0.6109,-0.9889) (3.8397,0.6109,-1.0050) (3.9270,0.6109,-1.0070) (4.0143,0.6109,-0.9966) (4.1015,0.6109,-0.9765) (4.1888,0.6109,-0.9496) (4.2761,0.6109,-0.9194) (4.3633,0.6109,-0.8890) (4.4506,0.6109,-0.8616) (4.5379,0.6109,-0.8400) (4.6251,0.6109,-0.8262) (4.7124,0.6109,-0.8214) (4.7997,0.6109,-0.8262) (4.8869,0.6109,-0.8400) (4.9742,0.6109,-0.8616) (5.0615,0.6109,-0.8890) (5.1487,0.6109,-0.9194) (5.2360,0.6109,-0.9496) (5.3233,0.6109,-0.9765) (5.4105,0.6109,-0.9966) (5.4978,0.6109,-1.0070) (5.5851,0.6109,-1.0050) (5.6723,0.6109,-0.9889) (5.7596,0.6109,-0.9574) (5.8469,0.6109,-0.9105) (5.9341,0.6109,-0.8489) (6.0214,0.6109,-0.7743) (6.1087,0.6109,-0.6891) (6.1959,0.6109,-0.5964) (6.2832,0.6109,-0.5000)
(0.0000,0.6981,-0.5000) (0.0873,0.6981,-0.4277) (0.1745,0.6981,-0.3586) (0.2618,0.6981,-0.2955) (0.3491,0.6981,-0.2412) (0.4363,0.6981,-0.1974) (0.5236,0.6981,-0.1657) (0.6109,0.6981,-0.1466) (0.6981,0.6981,-0.1398) (0.7854,0.6981,-0.1445) (0.8727,0.6981,-0.1590) (0.9599,0.6981,-0.1811) (1.0472,0.6981,-0.2083) (1.1345,0.6981,-0.2376) (1.2217,0.6981,-0.2664) (1.3090,0.6981,-0.2919) (1.3963,0.6981,-0.3119) (1.4835,0.6981,-0.3246) (1.5708,0.6981,-0.3290) (1.6581,0.6981,-0.3246) (1.7453,0.6981,-0.3119) (1.8326,0.6981,-0.2919) (1.9199,0.6981,-0.2664) (2.0071,0.6981,-0.2376) (2.0944,0.6981,-0.2083) (2.1817,0.6981,-0.1811) (2.2689,0.6981,-0.1590) (2.3562,0.6981,-0.1445) (2.4435,0.6981,-0.1398) (2.5307,0.6981,-0.1466) (2.6180,0.6981,-0.1657) (2.7053,0.6981,-0.1974) (2.7925,0.6981,-0.2412) (2.8798,0.6981,-0.2955) (2.9671,0.6981,-0.3586) (3.0543,0.6981,-0.4277) (3.1416,0.6981,-0.5000) (3.2289,0.6981,-0.5723) (3.3161,0.6981,-0.6414) (3.4034,0.6981,-0.7045) (3.4907,0.6981,-0.7588) (3.5779,0.6981,-0.8026) (3.6652,0.6981,-0.8343) (3.7525,0.6981,-0.8534) (3.8397,0.6981,-0.8602) (3.9270,0.6981,-0.8555) (4.0143,0.6981,-0.8410) (4.1015,0.6981,-0.8189) (4.1888,0.6981,-0.7917) (4.2761,0.6981,-0.7624) (4.3633,0.6981,-0.7336) (4.4506,0.6981,-0.7081) (4.5379,0.6981,-0.6881) (4.6251,0.6981,-0.6754) (4.7124,0.6981,-0.6710) (4.7997,0.6981,-0.6754) (4.8869,0.6981,-0.6881) (4.9742,0.6981,-0.7081) (5.0615,0.6981,-0.7336) (5.1487,0.6981,-0.7624) (5.2360,0.6981,-0.7917) (5.3233,0.6981,-0.8189) (5.4105,0.6981,-0.8410) (5.4978,0.6981,-0.8555) (5.5851,0.6981,-0.8602) (5.6723,0.6981,-0.8534) (5.7596,0.6981,-0.8343) (5.8469,0.6981,-0.8026) (5.9341,0.6981,-0.7588) (6.0214,0.6981,-0.7045) (6.1087,0.6981,-0.6414) (6.1959,0.6981,-0.5723) (6.2832,0.6981,-0.5000)
(0.0000,0.7854,-0.5000) (0.0873,0.7854,-0.4568) (0.1745,0.7854,-0.4158) (0.2618,0.7854,-0.3793) (0.3491,0.7854,-0.3490) (0.4363,0.7854,-0.3264) (0.5236,0.7854,-0.3125) (0.6109,0.7854,-0.3076) (0.6981,0.7854,-0.3114) (0.7854,0.7854,-0.3232) (0.8727,0.7854,-0.3417) (0.9599,0.7854,-0.3653) (1.0472,0.7854,-0.3917) (1.1345,0.7854,-0.4191) (1.2217,0.7854,-0.4450) (1.3090,0.7854,-0.4676) (1.3963,0.7854,-0.4852) (1.4835,0.7854,-0.4962) (1.5708,0.7854,-0.5000) (1.6581,0.7854,-0.4962) (1.7453,0.7854,-0.4852) (1.8326,0.7854,-0.4676) (1.9199,0.7854,-0.4450) (2.0071,0.7854,-0.4191) (2.0944,0.7854,-0.3917) (2.1817,0.7854,-0.3653) (2.2689,0.7854,-0.3417) (2.3562,0.7854,-0.3232) (2.4435,0.7854,-0.3114) (2.5307,0.7854,-0.3076) (2.6180,0.7854,-0.3125) (2.7053,0.7854,-0.3264) (2.7925,0.7854,-0.3490) (2.8798,0.7854,-0.3793) (2.9671,0.7854,-0.4158) (3.0543,0.7854,-0.4568) (3.1416,0.7854,-0.5000) (3.2289,0.7854,-0.5432) (3.3161,0.7854,-0.5842) (3.4034,0.7854,-0.6207) (3.4907,0.7854,-0.6510) (3.5779,0.7854,-0.6736) (3.6652,0.7854,-0.6875) (3.7525,0.7854,-0.6924) (3.8397,0.7854,-0.6886) (3.9270,0.7854,-0.6768) (4.0143,0.7854,-0.6583) (4.1015,0.7854,-0.6347) (4.1888,0.7854,-0.6083) (4.2761,0.7854,-0.5809) (4.3633,0.7854,-0.5550) (4.4506,0.7854,-0.5324) (4.5379,0.7854,-0.5148) (4.6251,0.7854,-0.5038) (4.7124,0.7854,-0.5000) (4.7997,0.7854,-0.5038) (4.8869,0.7854,-0.5148) (4.9742,0.7854,-0.5324) (5.0615,0.7854,-0.5550) (5.1487,0.7854,-0.5809) (5.2360,0.7854,-0.6083) (5.3233,0.7854,-0.6347) (5.4105,0.7854,-0.6583) (5.4978,0.7854,-0.6768) (5.5851,0.7854,-0.6886) (5.6723,0.7854,-0.6924) (5.7596,0.7854,-0.6875) (5.8469,0.7854,-0.6736) (5.9341,0.7854,-0.6510) (6.0214,0.7854,-0.6207) (6.1087,0.7854,-0.5842) (6.1959,0.7854,-0.5432) (6.2832,0.7854,-0.5000)
(0.0000,0.8727,-0.5000) (0.0873,0.8727,-0.4871) (0.1745,0.8727,-0.4756) (0.2618,0.8727,-0.4667) (0.3491,0.8727,-0.4614) (0.4363,0.8727,-0.4607) (0.5236,0.8727,-0.4650) (0.6109,0.8727,-0.4744) (0.6981,0.8727,-0.4887) (0.7854,0.8727,-0.5073) (0.8727,0.8727,-0.5293) (0.9599,0.8727,-0.5535) (1.0472,0.8727,-0.5785) (1.1345,0.8727,-0.6030) (1.2217,0.8727,-0.6254) (1.3090,0.8727,-0.6444) (1.3963,0.8727,-0.6589) (1.4835,0.8727,-0.6679) (1.5708,0.8727,-0.6710) (1.6581,0.8727,-0.6679) (1.7453,0.8727,-0.6589) (1.8326,0.8727,-0.6444) (1.9199,0.8727,-0.6254) (2.0071,0.8727,-0.6030) (2.0944,0.8727,-0.5785) (2.1817,0.8727,-0.5535) (2.2689,0.8727,-0.5293) (2.3562,0.8727,-0.5073) (2.4435,0.8727,-0.4887) (2.5307,0.8727,-0.4744) (2.6180,0.8727,-0.4650) (2.7053,0.8727,-0.4607) (2.7925,0.8727,-0.4614) (2.8798,0.8727,-0.4667) (2.9671,0.8727,-0.4756) (3.0543,0.8727,-0.4871) (3.1416,0.8727,-0.5000) (3.2289,0.8727,-0.5129) (3.3161,0.8727,-0.5244) (3.4034,0.8727,-0.5333) (3.4907,0.8727,-0.5386) (3.5779,0.8727,-0.5393) (3.6652,0.8727,-0.5350) (3.7525,0.8727,-0.5256) (3.8397,0.8727,-0.5113) (3.9270,0.8727,-0.4927) (4.0143,0.8727,-0.4707) (4.1015,0.8727,-0.4465) (4.1888,0.8727,-0.4215) (4.2761,0.8727,-0.3970) (4.3633,0.8727,-0.3746) (4.4506,0.8727,-0.3556) (4.5379,0.8727,-0.3411) (4.6251,0.8727,-0.3321) (4.7124,0.8727,-0.3290) (4.7997,0.8727,-0.3321) (4.8869,0.8727,-0.3411) (4.9742,0.8727,-0.3556) (5.0615,0.8727,-0.3746) (5.1487,0.8727,-0.3970) (5.2360,0.8727,-0.4215) (5.3233,0.8727,-0.4465) (5.4105,0.8727,-0.4707) (5.4978,0.8727,-0.4927) (5.5851,0.8727,-0.5113) (5.6723,0.8727,-0.5256) (5.7596,0.8727,-0.5350) (5.8469,0.8727,-0.5393) (5.9341,0.8727,-0.5386) (6.0214,0.8727,-0.5333) (6.1087,0.8727,-0.5244) (6.1959,0.8727,-0.5129) (6.2832,0.8727,-0.5000)
(0.0000,0.9599,-0.5000) (0.0873,0.9599,-0.5152) (0.1745,0.9599,-0.5308) (0.2618,0.9599,-0.5473) (0.3491,0.9599,-0.5651) (0.4363,0.9599,-0.5843) (0.5236,0.9599,-0.6050) (0.6109,0.9599,-0.6272) (0.6981,0.9599,-0.6506) (0.7854,0.9599,-0.6748) (0.8727,0.9599,-0.6992) (0.9599,0.9599,-0.7233) (1.0472,0.9599,-0.7462) (1.1345,0.9599,-0.7673) (1.2217,0.9599,-0.7857) (1.3090,0.9599,-0.8008) (1.3963,0.9599,-0.8121) (1.4835,0.9599,-0.8190) (1.5708,0.9599,-0.8214) (1.6581,0.9599,-0.8190) (1.7453,0.9599,-0.8121) (1.8326,0.9599,-0.8008) (1.9199,0.9599,-0.7857) (2.0071,0.9599,-0.7673) (2.0944,0.9599,-0.7462) (2.1817,0.9599,-0.7233) (2.2689,0.9599,-0.6992) (2.3562,0.9599,-0.6748) (2.4435,0.9599,-0.6506) (2.5307,0.9599,-0.6272) (2.6180,0.9599,-0.6050) (2.7053,0.9599,-0.5843) (2.7925,0.9599,-0.5651) (2.8798,0.9599,-0.5473) (2.9671,0.9599,-0.5308) (3.0543,0.9599,-0.5152) (3.1416,0.9599,-0.5000) (3.2289,0.9599,-0.4848) (3.3161,0.9599,-0.4692) (3.4034,0.9599,-0.4527) (3.4907,0.9599,-0.4349) (3.5779,0.9599,-0.4157) (3.6652,0.9599,-0.3950) (3.7525,0.9599,-0.3728) (3.8397,0.9599,-0.3494) (3.9270,0.9599,-0.3252) (4.0143,0.9599,-0.3008) (4.1015,0.9599,-0.2767) (4.1888,0.9599,-0.2538) (4.2761,0.9599,-0.2327) (4.3633,0.9599,-0.2143) (4.4506,0.9599,-0.1992) (4.5379,0.9599,-0.1879) (4.6251,0.9599,-0.1810) (4.7124,0.9599,-0.1786) (4.7997,0.9599,-0.1810) (4.8869,0.9599,-0.1879) (4.9742,0.9599,-0.1992) (5.0615,0.9599,-0.2143) (5.1487,0.9599,-0.2327) (5.2360,0.9599,-0.2538) (5.3233,0.9599,-0.2767) (5.4105,0.9599,-0.3008) (5.4978,0.9599,-0.3252) (5.5851,0.9599,-0.3494) (5.6723,0.9599,-0.3728) (5.7596,0.9599,-0.3950) (5.8469,0.9599,-0.4157) (5.9341,0.9599,-0.4349) (6.0214,0.9599,-0.4527) (6.1087,0.9599,-0.4692) (6.1959,0.9599,-0.4848) (6.2832,0.9599,-0.5000)
(0.0000,1.0472,-0.5000) (0.0873,1.0472,-0.5377) (0.1745,1.0472,-0.5752) (0.2618,1.0472,-0.6121) (0.3491,1.0472,-0.6481) (0.4363,1.0472,-0.6830) (0.5236,1.0472,-0.7165) (0.6109,1.0472,-0.7484) (0.6981,1.0472,-0.7783) (0.7854,1.0472,-0.8062) (0.8727,1.0472,-0.8317) (0.9599,1.0472,-0.8547) (1.0472,1.0472,-0.8750) (1.1345,1.0472,-0.8924) (1.2217,1.0472,-0.9069) (1.3090,1.0472,-0.9183) (1.3963,1.0472,-0.9264) (1.4835,1.0472,-0.9314) (1.5708,1.0472,-0.9330) (1.6581,1.0472,-0.9314) (1.7453,1.0472,-0.9264) (1.8326,1.0472,-0.9183) (1.9199,1.0472,-0.9069) (2.0071,1.0472,-0.8924) (2.0944,1.0472,-0.8750) (2.1817,1.0472,-0.8547) (2.2689,1.0472,-0.8317) (2.3562,1.0472,-0.8062) (2.4435,1.0472,-0.7783) (2.5307,1.0472,-0.7484) (2.6180,1.0472,-0.7165) (2.7053,1.0472,-0.6830) (2.7925,1.0472,-0.6481) (2.8798,1.0472,-0.6121) (2.9671,1.0472,-0.5752) (3.0543,1.0472,-0.5377) (3.1416,1.0472,-0.5000) (3.2289,1.0472,-0.4623) (3.3161,1.0472,-0.4248) (3.4034,1.0472,-0.3879) (3.4907,1.0472,-0.3519) (3.5779,1.0472,-0.3170) (3.6652,1.0472,-0.2835) (3.7525,1.0472,-0.2516) (3.8397,1.0472,-0.2217) (3.9270,1.0472,-0.1938) (4.0143,1.0472,-0.1683) (4.1015,1.0472,-0.1453) (4.1888,1.0472,-0.1250) (4.2761,1.0472,-0.1076) (4.3633,1.0472,-0.0931) (4.4506,1.0472,-0.0817) (4.5379,1.0472,-0.0736) (4.6251,1.0472,-0.0686) (4.7124,1.0472,-0.0670) (4.7997,1.0472,-0.0686) (4.8869,1.0472,-0.0736) (4.9742,1.0472,-0.0817) (5.0615,1.0472,-0.0931) (5.1487,1.0472,-0.1076) (5.2360,1.0472,-0.1250) (5.3233,1.0472,-0.1453) (5.4105,1.0472,-0.1683) (5.4978,1.0472,-0.1938) (5.5851,1.0472,-0.2217) (5.6723,1.0472,-0.2516) (5.7596,1.0472,-0.2835) (5.8469,1.0472,-0.3170) (5.9341,1.0472,-0.3519) (6.0214,1.0472,-0.3879) (6.1087,1.0472,-0.4248) (6.1959,1.0472,-0.4623) (6.2832,1.0472,-0.5000)
(0.0000,1.1345,-0.5000) (0.0873,1.1345,-0.5524) (0.1745,1.1345,-0.6039) (0.2618,1.1345,-0.6539) (0.3491,1.1345,-0.7014) (0.4363,1.1345,-0.7461) (0.5236,1.1345,-0.7872) (0.6109,1.1345,-0.8245) (0.6981,1.1345,-0.8578) (0.7854,1.1345,-0.8869) (0.8727,1.1345,-0.9118) (0.9599,1.1345,-0.9328) (1.0472,1.1345,-0.9501) (1.1345,1.1345,-0.9640) (1.2217,1.1345,-0.9747) (1.3090,1.1345,-0.9827) (1.3963,1.1345,-0.9882) (1.4835,1.1345,-0.9914) (1.5708,1.1345,-0.9924) (1.6581,1.1345,-0.9914) (1.7453,1.1345,-0.9882) (1.8326,1.1345,-0.9827) (1.9199,1.1345,-0.9747) (2.0071,1.1345,-0.9640) (2.0944,1.1345,-0.9501) (2.1817,1.1345,-0.9328) (2.2689,1.1345,-0.9118) (2.3562,1.1345,-0.8869) (2.4435,1.1345,-0.8578) (2.5307,1.1345,-0.8245) (2.6180,1.1345,-0.7872) (2.7053,1.1345,-0.7461) (2.7925,1.1345,-0.7014) (2.8798,1.1345,-0.6539) (2.9671,1.1345,-0.6039) (3.0543,1.1345,-0.5524) (3.1416,1.1345,-0.5000) (3.2289,1.1345,-0.4476) (3.3161,1.1345,-0.3961) (3.4034,1.1345,-0.3461) (3.4907,1.1345,-0.2986) (3.5779,1.1345,-0.2539) (3.6652,1.1345,-0.2128) (3.7525,1.1345,-0.1755) (3.8397,1.1345,-0.1422) (3.9270,1.1345,-0.1131) (4.0143,1.1345,-0.0882) (4.1015,1.1345,-0.0672) (4.1888,1.1345,-0.0499) (4.2761,1.1345,-0.0360) (4.3633,1.1345,-0.0253) (4.4506,1.1345,-0.0173) (4.5379,1.1345,-0.0118) (4.6251,1.1345,-0.0086) (4.7124,1.1345,-0.0076) (4.7997,1.1345,-0.0086) (4.8869,1.1345,-0.0118) (4.9742,1.1345,-0.0173) (5.0615,1.1345,-0.0253) (5.1487,1.1345,-0.0360) (5.2360,1.1345,-0.0499) (5.3233,1.1345,-0.0672) (5.4105,1.1345,-0.0882) (5.4978,1.1345,-0.1131) (5.5851,1.1345,-0.1422) (5.6723,1.1345,-0.1755) (5.7596,1.1345,-0.2128) (5.8469,1.1345,-0.2539) (5.9341,1.1345,-0.2986) (6.0214,1.1345,-0.3461) (6.1087,1.1345,-0.3961) (6.1959,1.1345,-0.4476) (6.2832,1.1345,-0.5000)
(0.0000,1.2217,-0.5000) (0.0873,1.2217,-0.5577) (0.1745,1.2217,-0.6143) (0.2618,1.2217,-0.6687) (0.3491,1.2217,-0.7201) (0.4363,1.2217,-0.7675) (0.5236,1.2217,-0.8103) (0.6109,1.2217,-0.8482) (0.6981,1.2217,-0.8810) (0.7854,1.2217,-0.9086) (0.8727,1.2217,-0.9313) (0.9599,1.2217,-0.9494) (1.0472,1.2217,-0.9635) (1.1345,1.2217,-0.9740) (1.2217,1.2217,-0.9815) (1.3090,1.2217,-0.9867) (1.3963,1.2217,-0.9900) (1.4835,1.2217,-0.9918) (1.5708,1.2217,-0.9924) (1.6581,1.2217,-0.9918) (1.7453,1.2217,-0.9900) (1.8326,1.2217,-0.9867) (1.9199,1.2217,-0.9815) (2.0071,1.2217,-0.9740) (2.0944,1.2217,-0.9635) (2.1817,1.2217,-0.9494) (2.2689,1.2217,-0.9313) (2.3562,1.2217,-0.9086) (2.4435,1.2217,-0.8810) (2.5307,1.2217,-0.8482) (2.6180,1.2217,-0.8103) (2.7053,1.2217,-0.7675) (2.7925,1.2217,-0.7201) (2.8798,1.2217,-0.6687) (2.9671,1.2217,-0.6143) (3.0543,1.2217,-0.5577) (3.1416,1.2217,-0.5000) (3.2289,1.2217,-0.4423) (3.3161,1.2217,-0.3857) (3.4034,1.2217,-0.3313) (3.4907,1.2217,-0.2799) (3.5779,1.2217,-0.2325) (3.6652,1.2217,-0.1897) (3.7525,1.2217,-0.1518) (3.8397,1.2217,-0.1190) (3.9270,1.2217,-0.0914) (4.0143,1.2217,-0.0687) (4.1015,1.2217,-0.0506) (4.1888,1.2217,-0.0365) (4.2761,1.2217,-0.0260) (4.3633,1.2217,-0.0185) (4.4506,1.2217,-0.0133) (4.5379,1.2217,-0.0100) (4.6251,1.2217,-0.0082) (4.7124,1.2217,-0.0076) (4.7997,1.2217,-0.0082) (4.8869,1.2217,-0.0100) (4.9742,1.2217,-0.0133) (5.0615,1.2217,-0.0185) (5.1487,1.2217,-0.0260) (5.2360,1.2217,-0.0365) (5.3233,1.2217,-0.0506) (5.4105,1.2217,-0.0687) (5.4978,1.2217,-0.0914) (5.5851,1.2217,-0.1190) (5.6723,1.2217,-0.1518) (5.7596,1.2217,-0.1897) (5.8469,1.2217,-0.2325) (5.9341,1.2217,-0.2799) (6.0214,1.2217,-0.3313) (6.1087,1.2217,-0.3857) (6.1959,1.2217,-0.4423) (6.2832,1.2217,-0.5000)
(0.0000,1.3090,-0.5000) (0.0873,1.3090,-0.5536) (0.1745,1.3090,-0.6060) (0.2618,1.3090,-0.6563) (0.3491,1.3090,-0.7034) (0.4363,1.3090,-0.7465) (0.5236,1.3090,-0.7851) (0.6109,1.3090,-0.8188) (0.6981,1.3090,-0.8474) (0.7854,1.3090,-0.8709) (0.8727,1.3090,-0.8896) (0.9599,1.3090,-0.9040) (1.0472,1.3090,-0.9146) (1.1345,1.3090,-0.9221) (1.2217,1.3090,-0.9270) (1.3090,1.3090,-0.9301) (1.3963,1.3090,-0.9319) (1.4835,1.3090,-0.9327) (1.5708,1.3090,-0.9330) (1.6581,1.3090,-0.9327) (1.7453,1.3090,-0.9319) (1.8326,1.3090,-0.9301) (1.9199,1.3090,-0.9270) (2.0071,1.3090,-0.9221) (2.0944,1.3090,-0.9146) (2.1817,1.3090,-0.9040) (2.2689,1.3090,-0.8896) (2.3562,1.3090,-0.8709) (2.4435,1.3090,-0.8474) (2.5307,1.3090,-0.8188) (2.6180,1.3090,-0.7851) (2.7053,1.3090,-0.7465) (2.7925,1.3090,-0.7034) (2.8798,1.3090,-0.6563) (2.9671,1.3090,-0.6060) (3.0543,1.3090,-0.5536) (3.1416,1.3090,-0.5000) (3.2289,1.3090,-0.4464) (3.3161,1.3090,-0.3940) (3.4034,1.3090,-0.3437) (3.4907,1.3090,-0.2966) (3.5779,1.3090,-0.2535) (3.6652,1.3090,-0.2149) (3.7525,1.3090,-0.1812) (3.8397,1.3090,-0.1526) (3.9270,1.3090,-0.1291) (4.0143,1.3090,-0.1104) (4.1015,1.3090,-0.0960) (4.1888,1.3090,-0.0854) (4.2761,1.3090,-0.0779) (4.3633,1.3090,-0.0730) (4.4506,1.3090,-0.0699) (4.5379,1.3090,-0.0681) (4.6251,1.3090,-0.0673) (4.7124,1.3090,-0.0670) (4.7997,1.3090,-0.0673) (4.8869,1.3090,-0.0681) (4.9742,1.3090,-0.0699) (5.0615,1.3090,-0.0730) (5.1487,1.3090,-0.0779) (5.2360,1.3090,-0.0854) (5.3233,1.3090,-0.0960) (5.4105,1.3090,-0.1104) (5.4978,1.3090,-0.1291) (5.5851,1.3090,-0.1526) (5.6723,1.3090,-0.1812) (5.7596,1.3090,-0.2149) (5.8469,1.3090,-0.2535) (5.9341,1.3090,-0.2966) (6.0214,1.3090,-0.3437) (6.1087,1.3090,-0.3940) (6.1959,1.3090,-0.4464) (6.2832,1.3090,-0.5000)
(0.0000,1.3963,-0.5000) (0.0873,1.3963,-0.5410) (0.1745,1.3963,-0.5811) (0.2618,1.3963,-0.6195) (0.3491,1.3963,-0.6553) (0.4363,1.3963,-0.6880) (0.5236,1.3963,-0.7171) (0.6109,1.3963,-0.7422) (0.6981,1.3963,-0.7633) (0.7854,1.3963,-0.7804) (0.8727,1.3963,-0.7938) (0.9599,1.3963,-0.8038) (1.0472,1.3963,-0.8109) (1.1345,1.3963,-0.8156) (1.2217,1.3963,-0.8185) (1.3090,1.3963,-0.8202) (1.3963,1.3963,-0.8210) (1.4835,1.3963,-0.8213) (1.5708,1.3963,-0.8214) (1.6581,1.3963,-0.8213) (1.7453,1.3963,-0.8210) (1.8326,1.3963,-0.8202) (1.9199,1.3963,-0.8185) (2.0071,1.3963,-0.8156) (2.0944,1.3963,-0.8109) (2.1817,1.3963,-0.8038) (2.2689,1.3963,-0.7938) (2.3562,1.3963,-0.7804) (2.4435,1.3963,-0.7633) (2.5307,1.3963,-0.7422) (2.6180,1.3963,-0.7171) (2.7053,1.3963,-0.6880) (2.7925,1.3963,-0.6553) (2.8798,1.3963,-0.6195) (2.9671,1.3963,-0.5811) (3.0543,1.3963,-0.5410) (3.1416,1.3963,-0.5000) (3.2289,1.3963,-0.4590) (3.3161,1.3963,-0.4189) (3.4034,1.3963,-0.3805) (3.4907,1.3963,-0.3447) (3.5779,1.3963,-0.3120) (3.6652,1.3963,-0.2829) (3.7525,1.3963,-0.2578) (3.8397,1.3963,-0.2367) (3.9270,1.3963,-0.2196) (4.0143,1.3963,-0.2062) (4.1015,1.3963,-0.1962) (4.1888,1.3963,-0.1891) (4.2761,1.3963,-0.1844) (4.3633,1.3963,-0.1815) (4.4506,1.3963,-0.1798) (4.5379,1.3963,-0.1790) (4.6251,1.3963,-0.1787) (4.7124,1.3963,-0.1786) (4.7997,1.3963,-0.1787) (4.8869,1.3963,-0.1790) (4.9742,1.3963,-0.1798) (5.0615,1.3963,-0.1815) (5.1487,1.3963,-0.1844) (5.2360,1.3963,-0.1891) (5.3233,1.3963,-0.1962) (5.4105,1.3963,-0.2062) (5.4978,1.3963,-0.2196) (5.5851,1.3963,-0.2367) (5.6723,1.3963,-0.2578) (5.7596,1.3963,-0.2829) (5.8469,1.3963,-0.3120) (5.9341,1.3963,-0.3447) (6.0214,1.3963,-0.3805) (6.1087,1.3963,-0.4189) (6.1959,1.3963,-0.4590) (6.2832,1.3963,-0.5000)
(0.0000,1.4835,-0.5000) (0.0873,1.4835,-0.5222) (0.1745,1.4835,-0.5439) (0.2618,1.4835,-0.5646) (0.3491,1.4835,-0.5839) (0.4363,1.4835,-0.6015) (0.5236,1.4835,-0.6171) (0.6109,1.4835,-0.6305) (0.6981,1.4835,-0.6417) (0.7854,1.4835,-0.6507) (0.8727,1.4835,-0.6576) (0.9599,1.4835,-0.6628) (1.0472,1.4835,-0.6663) (1.1345,1.4835,-0.6686) (1.2217,1.4835,-0.6700) (1.3090,1.4835,-0.6706) (1.3963,1.4835,-0.6709) (1.4835,1.4835,-0.6710) (1.5708,1.4835,-0.6710) (1.6581,1.4835,-0.6710) (1.7453,1.4835,-0.6709) (1.8326,1.4835,-0.6706) (1.9199,1.4835,-0.6700) (2.0071,1.4835,-0.6686) (2.0944,1.4835,-0.6663) (2.1817,1.4835,-0.6628) (2.2689,1.4835,-0.6576) (2.3562,1.4835,-0.6507) (2.4435,1.4835,-0.6417) (2.5307,1.4835,-0.6305) (2.6180,1.4835,-0.6171) (2.7053,1.4835,-0.6015) (2.7925,1.4835,-0.5839) (2.8798,1.4835,-0.5646) (2.9671,1.4835,-0.5439) (3.0543,1.4835,-0.5222) (3.1416,1.4835,-0.5000) (3.2289,1.4835,-0.4778) (3.3161,1.4835,-0.4561) (3.4034,1.4835,-0.4354) (3.4907,1.4835,-0.4161) (3.5779,1.4835,-0.3985) (3.6652,1.4835,-0.3829) (3.7525,1.4835,-0.3695) (3.8397,1.4835,-0.3583) (3.9270,1.4835,-0.3493) (4.0143,1.4835,-0.3424) (4.1015,1.4835,-0.3372) (4.1888,1.4835,-0.3337) (4.2761,1.4835,-0.3314) (4.3633,1.4835,-0.3300) (4.4506,1.4835,-0.3294) (4.5379,1.4835,-0.3291) (4.6251,1.4835,-0.3290) (4.7124,1.4835,-0.3290) (4.7997,1.4835,-0.3290) (4.8869,1.4835,-0.3291) (4.9742,1.4835,-0.3294) (5.0615,1.4835,-0.3300) (5.1487,1.4835,-0.3314) (5.2360,1.4835,-0.3337) (5.3233,1.4835,-0.3372) (5.4105,1.4835,-0.3424) (5.4978,1.4835,-0.3493) (5.5851,1.4835,-0.3583) (5.6723,1.4835,-0.3695) (5.7596,1.4835,-0.3829) (5.8469,1.4835,-0.3985) (5.9341,1.4835,-0.4161) (6.0214,1.4835,-0.4354) (6.1087,1.4835,-0.4561) (6.1959,1.4835,-0.4778) (6.2832,1.4835,-0.5000)
(0.0000,1.5708,-0.5000) (0.0873,1.5708,-0.5000) (0.1745,1.5708,-0.5000) (0.2618,1.5708,-0.5000) (0.3491,1.5708,-0.5000) (0.4363,1.5708,-0.5000) (0.5236,1.5708,-0.5000) (0.6109,1.5708,-0.5000) (0.6981,1.5708,-0.5000) (0.7854,1.5708,-0.5000) (0.8727,1.5708,-0.5000) (0.9599,1.5708,-0.5000) (1.0472,1.5708,-0.5000) (1.1345,1.5708,-0.5000) (1.2217,1.5708,-0.5000) (1.3090,1.5708,-0.5000) (1.3963,1.5708,-0.5000) (1.4835,1.5708,-0.5000) (1.5708,1.5708,-0.5000) (1.6581,1.5708,-0.5000) (1.7453,1.5708,-0.5000) (1.8326,1.5708,-0.5000) (1.9199,1.5708,-0.5000) (2.0071,1.5708,-0.5000) (2.0944,1.5708,-0.5000) (2.1817,1.5708,-0.5000) (2.2689,1.5708,-0.5000) (2.3562,1.5708,-0.5000) (2.4435,1.5708,-0.5000) (2.5307,1.5708,-0.5000) (2.6180,1.5708,-0.5000) (2.7053,1.5708,-0.5000) (2.7925,1.5708,-0.5000) (2.8798,1.5708,-0.5000) (2.9671,1.5708,-0.5000) (3.0543,1.5708,-0.5000) (3.1416,1.5708,-0.5000) (3.2289,1.5708,-0.5000) (3.3161,1.5708,-0.5000) (3.4034,1.5708,-0.5000) (3.4907,1.5708,-0.5000) (3.5779,1.5708,-0.5000) (3.6652,1.5708,-0.5000) (3.7525,1.5708,-0.5000) (3.8397,1.5708,-0.5000) (3.9270,1.5708,-0.5000) (4.0143,1.5708,-0.5000) (4.1015,1.5708,-0.5000) (4.1888,1.5708,-0.5000) (4.2761,1.5708,-0.5000) (4.3633,1.5708,-0.5000) (4.4506,1.5708,-0.5000) (4.5379,1.5708,-0.5000) (4.6251,1.5708,-0.5000) (4.7124,1.5708,-0.5000) (4.7997,1.5708,-0.5000) (4.8869,1.5708,-0.5000) (4.9742,1.5708,-0.5000) (5.0615,1.5708,-0.5000) (5.1487,1.5708,-0.5000) (5.2360,1.5708,-0.5000) (5.3233,1.5708,-0.5000) (5.4105,1.5708,-0.5000) (5.4978,1.5708,-0.5000) (5.5851,1.5708,-0.5000) (5.6723,1.5708,-0.5000) (5.7596,1.5708,-0.5000) (5.8469,1.5708,-0.5000) (5.9341,1.5708,-0.5000) (6.0214,1.5708,-0.5000) (6.1087,1.5708,-0.5000) (6.1959,1.5708,-0.5000) (6.2832,1.5708,-0.5000)
(0.0000,1.6581,-0.5000) (0.0873,1.6581,-0.4778) (0.1745,1.6581,-0.4561) (0.2618,1.6581,-0.4354) (0.3491,1.6581,-0.4161) (0.4363,1.6581,-0.3985) (0.5236,1.6581,-0.3829) (0.6109,1.6581,-0.3695) (0.6981,1.6581,-0.3583) (0.7854,1.6581,-0.3493) (0.8727,1.6581,-0.3424) (0.9599,1.6581,-0.3372) (1.0472,1.6581,-0.3337) (1.1345,1.6581,-0.3314) (1.2217,1.6581,-0.3300) (1.3090,1.6581,-0.3294) (1.3963,1.6581,-0.3291) (1.4835,1.6581,-0.3290) (1.5708,1.6581,-0.3290) (1.6581,1.6581,-0.3290) (1.7453,1.6581,-0.3291) (1.8326,1.6581,-0.3294) (1.9199,1.6581,-0.3300) (2.0071,1.6581,-0.3314) (2.0944,1.6581,-0.3337) (2.1817,1.6581,-0.3372) (2.2689,1.6581,-0.3424) (2.3562,1.6581,-0.3493) (2.4435,1.6581,-0.3583) (2.5307,1.6581,-0.3695) (2.6180,1.6581,-0.3829) (2.7053,1.6581,-0.3985) (2.7925,1.6581,-0.4161) (2.8798,1.6581,-0.4354) (2.9671,1.6581,-0.4561) (3.0543,1.6581,-0.4778) (3.1416,1.6581,-0.5000) (3.2289,1.6581,-0.5222) (3.3161,1.6581,-0.5439) (3.4034,1.6581,-0.5646) (3.4907,1.6581,-0.5839) (3.5779,1.6581,-0.6015) (3.6652,1.6581,-0.6171) (3.7525,1.6581,-0.6305) (3.8397,1.6581,-0.6417) (3.9270,1.6581,-0.6507) (4.0143,1.6581,-0.6576) (4.1015,1.6581,-0.6628) (4.1888,1.6581,-0.6663) (4.2761,1.6581,-0.6686) (4.3633,1.6581,-0.6700) (4.4506,1.6581,-0.6706) (4.5379,1.6581,-0.6709) (4.6251,1.6581,-0.6710) (4.7124,1.6581,-0.6710) (4.7997,1.6581,-0.6710) (4.8869,1.6581,-0.6709) (4.9742,1.6581,-0.6706) (5.0615,1.6581,-0.6700) (5.1487,1.6581,-0.6686) (5.2360,1.6581,-0.6663) (5.3233,1.6581,-0.6628) (5.4105,1.6581,-0.6576) (5.4978,1.6581,-0.6507) (5.5851,1.6581,-0.6417) (5.6723,1.6581,-0.6305) (5.7596,1.6581,-0.6171) (5.8469,1.6581,-0.6015) (5.9341,1.6581,-0.5839) (6.0214,1.6581,-0.5646) (6.1087,1.6581,-0.5439) (6.1959,1.6581,-0.5222) (6.2832,1.6581,-0.5000)
(0.0000,1.7453,-0.5000) (0.0873,1.7453,-0.4590) (0.1745,1.7453,-0.4189) (0.2618,1.7453,-0.3805) (0.3491,1.7453,-0.3447) (0.4363,1.7453,-0.3120) (0.5236,1.7453,-0.2829) (0.6109,1.7453,-0.2578) (0.6981,1.7453,-0.2367) (0.7854,1.7453,-0.2196) (0.8727,1.7453,-0.2062) (0.9599,1.7453,-0.1962) (1.0472,1.7453,-0.1891) (1.1345,1.7453,-0.1844) (1.2217,1.7453,-0.1815) (1.3090,1.7453,-0.1798) (1.3963,1.7453,-0.1790) (1.4835,1.7453,-0.1787) (1.5708,1.7453,-0.1786) (1.6581,1.7453,-0.1787) (1.7453,1.7453,-0.1790) (1.8326,1.7453,-0.1798) (1.9199,1.7453,-0.1815) (2.0071,1.7453,-0.1844) (2.0944,1.7453,-0.1891) (2.1817,1.7453,-0.1962) (2.2689,1.7453,-0.2062) (2.3562,1.7453,-0.2196) (2.4435,1.7453,-0.2367) (2.5307,1.7453,-0.2578) (2.6180,1.7453,-0.2829) (2.7053,1.7453,-0.3120) (2.7925,1.7453,-0.3447) (2.8798,1.7453,-0.3805) (2.9671,1.7453,-0.4189) (3.0543,1.7453,-0.4590) (3.1416,1.7453,-0.5000) (3.2289,1.7453,-0.5410) (3.3161,1.7453,-0.5811) (3.4034,1.7453,-0.6195) (3.4907,1.7453,-0.6553) (3.5779,1.7453,-0.6880) (3.6652,1.7453,-0.7171) (3.7525,1.7453,-0.7422) (3.8397,1.7453,-0.7633) (3.9270,1.7453,-0.7804) (4.0143,1.7453,-0.7938) (4.1015,1.7453,-0.8038) (4.1888,1.7453,-0.8109) (4.2761,1.7453,-0.8156) (4.3633,1.7453,-0.8185) (4.4506,1.7453,-0.8202) (4.5379,1.7453,-0.8210) (4.6251,1.7453,-0.8213) (4.7124,1.7453,-0.8214) (4.7997,1.7453,-0.8213) (4.8869,1.7453,-0.8210) (4.9742,1.7453,-0.8202) (5.0615,1.7453,-0.8185) (5.1487,1.7453,-0.8156) (5.2360,1.7453,-0.8109) (5.3233,1.7453,-0.8038) (5.4105,1.7453,-0.7938) (5.4978,1.7453,-0.7804) (5.5851,1.7453,-0.7633) (5.6723,1.7453,-0.7422) (5.7596,1.7453,-0.7171) (5.8469,1.7453,-0.6880) (5.9341,1.7453,-0.6553) (6.0214,1.7453,-0.6195) (6.1087,1.7453,-0.5811) (6.1959,1.7453,-0.5410) (6.2832,1.7453,-0.5000)
(0.0000,1.8326,-0.5000) (0.0873,1.8326,-0.4464) (0.1745,1.8326,-0.3940) (0.2618,1.8326,-0.3437) (0.3491,1.8326,-0.2966) (0.4363,1.8326,-0.2535) (0.5236,1.8326,-0.2149) (0.6109,1.8326,-0.1812) (0.6981,1.8326,-0.1526) (0.7854,1.8326,-0.1291) (0.8727,1.8326,-0.1104) (0.9599,1.8326,-0.0960) (1.0472,1.8326,-0.0854) (1.1345,1.8326,-0.0779) (1.2217,1.8326,-0.0730) (1.3090,1.8326,-0.0699) (1.3963,1.8326,-0.0681) (1.4835,1.8326,-0.0673) (1.5708,1.8326,-0.0670) (1.6581,1.8326,-0.0673) (1.7453,1.8326,-0.0681) (1.8326,1.8326,-0.0699) (1.9199,1.8326,-0.0730) (2.0071,1.8326,-0.0779) (2.0944,1.8326,-0.0854) (2.1817,1.8326,-0.0960) (2.2689,1.8326,-0.1104) (2.3562,1.8326,-0.1291) (2.4435,1.8326,-0.1526) (2.5307,1.8326,-0.1812) (2.6180,1.8326,-0.2149) (2.7053,1.8326,-0.2535) (2.7925,1.8326,-0.2966) (2.8798,1.8326,-0.3437) (2.9671,1.8326,-0.3940) (3.0543,1.8326,-0.4464) (3.1416,1.8326,-0.5000) (3.2289,1.8326,-0.5536) (3.3161,1.8326,-0.6060) (3.4034,1.8326,-0.6563) (3.4907,1.8326,-0.7034) (3.5779,1.8326,-0.7465) (3.6652,1.8326,-0.7851) (3.7525,1.8326,-0.8188) (3.8397,1.8326,-0.8474) (3.9270,1.8326,-0.8709) (4.0143,1.8326,-0.8896) (4.1015,1.8326,-0.9040) (4.1888,1.8326,-0.9146) (4.2761,1.8326,-0.9221) (4.3633,1.8326,-0.9270) (4.4506,1.8326,-0.9301) (4.5379,1.8326,-0.9319) (4.6251,1.8326,-0.9327) (4.7124,1.8326,-0.9330) (4.7997,1.8326,-0.9327) (4.8869,1.8326,-0.9319) (4.9742,1.8326,-0.9301) (5.0615,1.8326,-0.9270) (5.1487,1.8326,-0.9221) (5.2360,1.8326,-0.9146) (5.3233,1.8326,-0.9040) (5.4105,1.8326,-0.8896) (5.4978,1.8326,-0.8709) (5.5851,1.8326,-0.8474) (5.6723,1.8326,-0.8188) (5.7596,1.8326,-0.7851) (5.8469,1.8326,-0.7465) (5.9341,1.8326,-0.7034) (6.0214,1.8326,-0.6563) (6.1087,1.8326,-0.6060) (6.1959,1.8326,-0.5536) (6.2832,1.8326,-0.5000)
(0.0000,1.9199,-0.5000) (0.0873,1.9199,-0.4423) (0.1745,1.9199,-0.3857) (0.2618,1.9199,-0.3313) (0.3491,1.9199,-0.2799) (0.4363,1.9199,-0.2325) (0.5236,1.9199,-0.1897) (0.6109,1.9199,-0.1518) (0.6981,1.9199,-0.1190) (0.7854,1.9199,-0.0914) (0.8727,1.9199,-0.0687) (0.9599,1.9199,-0.0506) (1.0472,1.9199,-0.0365) (1.1345,1.9199,-0.0260) (1.2217,1.9199,-0.0185) (1.3090,1.9199,-0.0133) (1.3963,1.9199,-0.0100) (1.4835,1.9199,-0.0082) (1.5708,1.9199,-0.0076) (1.6581,1.9199,-0.0082) (1.7453,1.9199,-0.0100) (1.8326,1.9199,-0.0133) (1.9199,1.9199,-0.0185) (2.0071,1.9199,-0.0260) (2.0944,1.9199,-0.0365) (2.1817,1.9199,-0.0506) (2.2689,1.9199,-0.0687) (2.3562,1.9199,-0.0914) (2.4435,1.9199,-0.1190) (2.5307,1.9199,-0.1518) (2.6180,1.9199,-0.1897) (2.7053,1.9199,-0.2325) (2.7925,1.9199,-0.2799) (2.8798,1.9199,-0.3313) (2.9671,1.9199,-0.3857) (3.0543,1.9199,-0.4423) (3.1416,1.9199,-0.5000) (3.2289,1.9199,-0.5577) (3.3161,1.9199,-0.6143) (3.4034,1.9199,-0.6687) (3.4907,1.9199,-0.7201) (3.5779,1.9199,-0.7675) (3.6652,1.9199,-0.8103) (3.7525,1.9199,-0.8482) (3.8397,1.9199,-0.8810) (3.9270,1.9199,-0.9086) (4.0143,1.9199,-0.9313) (4.1015,1.9199,-0.9494) (4.1888,1.9199,-0.9635) (4.2761,1.9199,-0.9740) (4.3633,1.9199,-0.9815) (4.4506,1.9199,-0.9867) (4.5379,1.9199,-0.9900) (4.6251,1.9199,-0.9918) (4.7124,1.9199,-0.9924) (4.7997,1.9199,-0.9918) (4.8869,1.9199,-0.9900) (4.9742,1.9199,-0.9867) (5.0615,1.9199,-0.9815) (5.1487,1.9199,-0.9740) (5.2360,1.9199,-0.9635) (5.3233,1.9199,-0.9494) (5.4105,1.9199,-0.9313) (5.4978,1.9199,-0.9086) (5.5851,1.9199,-0.8810) (5.6723,1.9199,-0.8482) (5.7596,1.9199,-0.8103) (5.8469,1.9199,-0.7675) (5.9341,1.9199,-0.7201) (6.0214,1.9199,-0.6687) (6.1087,1.9199,-0.6143) (6.1959,1.9199,-0.5577) (6.2832,1.9199,-0.5000)
(0.0000,2.0071,-0.5000) (0.0873,2.0071,-0.4476) (0.1745,2.0071,-0.3961) (0.2618,2.0071,-0.3461) (0.3491,2.0071,-0.2986) (0.4363,2.0071,-0.2539) (0.5236,2.0071,-0.2128) (0.6109,2.0071,-0.1755) (0.6981,2.0071,-0.1422) (0.7854,2.0071,-0.1131) (0.8727,2.0071,-0.0882) (0.9599,2.0071,-0.0672) (1.0472,2.0071,-0.0499) (1.1345,2.0071,-0.0360) (1.2217,2.0071,-0.0253) (1.3090,2.0071,-0.0173) (1.3963,2.0071,-0.0118) (1.4835,2.0071,-0.0086) (1.5708,2.0071,-0.0076) (1.6581,2.0071,-0.0086) (1.7453,2.0071,-0.0118) (1.8326,2.0071,-0.0173) (1.9199,2.0071,-0.0253) (2.0071,2.0071,-0.0360) (2.0944,2.0071,-0.0499) (2.1817,2.0071,-0.0672) (2.2689,2.0071,-0.0882) (2.3562,2.0071,-0.1131) (2.4435,2.0071,-0.1422) (2.5307,2.0071,-0.1755) (2.6180,2.0071,-0.2128) (2.7053,2.0071,-0.2539) (2.7925,2.0071,-0.2986) (2.8798,2.0071,-0.3461) (2.9671,2.0071,-0.3961) (3.0543,2.0071,-0.4476) (3.1416,2.0071,-0.5000) (3.2289,2.0071,-0.5524) (3.3161,2.0071,-0.6039) (3.4034,2.0071,-0.6539) (3.4907,2.0071,-0.7014) (3.5779,2.0071,-0.7461) (3.6652,2.0071,-0.7872) (3.7525,2.0071,-0.8245) (3.8397,2.0071,-0.8578) (3.9270,2.0071,-0.8869) (4.0143,2.0071,-0.9118) (4.1015,2.0071,-0.9328) (4.1888,2.0071,-0.9501) (4.2761,2.0071,-0.9640) (4.3633,2.0071,-0.9747) (4.4506,2.0071,-0.9827) (4.5379,2.0071,-0.9882) (4.6251,2.0071,-0.9914) (4.7124,2.0071,-0.9924) (4.7997,2.0071,-0.9914) (4.8869,2.0071,-0.9882) (4.9742,2.0071,-0.9827) (5.0615,2.0071,-0.9747) (5.1487,2.0071,-0.9640) (5.2360,2.0071,-0.9501) (5.3233,2.0071,-0.9328) (5.4105,2.0071,-0.9118) (5.4978,2.0071,-0.8869) (5.5851,2.0071,-0.8578) (5.6723,2.0071,-0.8245) (5.7596,2.0071,-0.7872) (5.8469,2.0071,-0.7461) (5.9341,2.0071,-0.7014) (6.0214,2.0071,-0.6539) (6.1087,2.0071,-0.6039) (6.1959,2.0071,-0.5524) (6.2832,2.0071,-0.5000)
(0.0000,2.0944,-0.5000) (0.0873,2.0944,-0.4623) (0.1745,2.0944,-0.4248) (0.2618,2.0944,-0.3879) (0.3491,2.0944,-0.3519) (0.4363,2.0944,-0.3170) (0.5236,2.0944,-0.2835) (0.6109,2.0944,-0.2516) (0.6981,2.0944,-0.2217) (0.7854,2.0944,-0.1938) (0.8727,2.0944,-0.1683) (0.9599,2.0944,-0.1453) (1.0472,2.0944,-0.1250) (1.1345,2.0944,-0.1076) (1.2217,2.0944,-0.0931) (1.3090,2.0944,-0.0817) (1.3963,2.0944,-0.0736) (1.4835,2.0944,-0.0686) (1.5708,2.0944,-0.0670) (1.6581,2.0944,-0.0686) (1.7453,2.0944,-0.0736) (1.8326,2.0944,-0.0817) (1.9199,2.0944,-0.0931) (2.0071,2.0944,-0.1076) (2.0944,2.0944,-0.1250) (2.1817,2.0944,-0.1453) (2.2689,2.0944,-0.1683) (2.3562,2.0944,-0.1938) (2.4435,2.0944,-0.2217) (2.5307,2.0944,-0.2516) (2.6180,2.0944,-0.2835) (2.7053,2.0944,-0.3170) (2.7925,2.0944,-0.3519) (2.8798,2.0944,-0.3879) (2.9671,2.0944,-0.4248) (3.0543,2.0944,-0.4623) (3.1416,2.0944,-0.5000) (3.2289,2.0944,-0.5377) (3.3161,2.0944,-0.5752) (3.4034,2.0944,-0.6121) (3.4907,2.0944,-0.6481) (3.5779,2.0944,-0.6830) (3.6652,2.0944,-0.7165) (3.7525,2.0944,-0.7484) (3.8397,2.0944,-0.7783) (3.9270,2.0944,-0.8062) (4.0143,2.0944,-0.8317) (4.1015,2.0944,-0.8547) (4.1888,2.0944,-0.8750) (4.2761,2.0944,-0.8924) (4.3633,2.0944,-0.9069) (4.4506,2.0944,-0.9183) (4.5379,2.0944,-0.9264) (4.6251,2.0944,-0.9314) (4.7124,2.0944,-0.9330) (4.7997,2.0944,-0.9314) (4.8869,2.0944,-0.9264) (4.9742,2.0944,-0.9183) (5.0615,2.0944,-0.9069) (5.1487,2.0944,-0.8924) (5.2360,2.0944,-0.8750) (5.3233,2.0944,-0.8547) (5.4105,2.0944,-0.8317) (5.4978,2.0944,-0.8062) (5.5851,2.0944,-0.7783) (5.6723,2.0944,-0.7484) (5.7596,2.0944,-0.7165) (5.8469,2.0944,-0.6830) (5.9341,2.0944,-0.6481) (6.0214,2.0944,-0.6121) (6.1087,2.0944,-0.5752) (6.1959,2.0944,-0.5377) (6.2832,2.0944,-0.5000)
(0.0000,2.1817,-0.5000) (0.0873,2.1817,-0.4848) (0.1745,2.1817,-0.4692) (0.2618,2.1817,-0.4527) (0.3491,2.1817,-0.4349) (0.4363,2.1817,-0.4157) (0.5236,2.1817,-0.3950) (0.6109,2.1817,-0.3728) (0.6981,2.1817,-0.3494) (0.7854,2.1817,-0.3252) (0.8727,2.1817,-0.3008) (0.9599,2.1817,-0.2767) (1.0472,2.1817,-0.2538) (1.1345,2.1817,-0.2327) (1.2217,2.1817,-0.2143) (1.3090,2.1817,-0.1992) (1.3963,2.1817,-0.1879) (1.4835,2.1817,-0.1810) (1.5708,2.1817,-0.1786) (1.6581,2.1817,-0.1810) (1.7453,2.1817,-0.1879) (1.8326,2.1817,-0.1992) (1.9199,2.1817,-0.2143) (2.0071,2.1817,-0.2327) (2.0944,2.1817,-0.2538) (2.1817,2.1817,-0.2767) (2.2689,2.1817,-0.3008) (2.3562,2.1817,-0.3252) (2.4435,2.1817,-0.3494) (2.5307,2.1817,-0.3728) (2.6180,2.1817,-0.3950) (2.7053,2.1817,-0.4157) (2.7925,2.1817,-0.4349) (2.8798,2.1817,-0.4527) (2.9671,2.1817,-0.4692) (3.0543,2.1817,-0.4848) (3.1416,2.1817,-0.5000) (3.2289,2.1817,-0.5152) (3.3161,2.1817,-0.5308) (3.4034,2.1817,-0.5473) (3.4907,2.1817,-0.5651) (3.5779,2.1817,-0.5843) (3.6652,2.1817,-0.6050) (3.7525,2.1817,-0.6272) (3.8397,2.1817,-0.6506) (3.9270,2.1817,-0.6748) (4.0143,2.1817,-0.6992) (4.1015,2.1817,-0.7233) (4.1888,2.1817,-0.7462) (4.2761,2.1817,-0.7673) (4.3633,2.1817,-0.7857) (4.4506,2.1817,-0.8008) (4.5379,2.1817,-0.8121) (4.6251,2.1817,-0.8190) (4.7124,2.1817,-0.8214) (4.7997,2.1817,-0.8190) (4.8869,2.1817,-0.8121) (4.9742,2.1817,-0.8008) (5.0615,2.1817,-0.7857) (5.1487,2.1817,-0.7673) (5.2360,2.1817,-0.7462) (5.3233,2.1817,-0.7233) (5.4105,2.1817,-0.6992) (5.4978,2.1817,-0.6748) (5.5851,2.1817,-0.6506) (5.6723,2.1817,-0.6272) (5.7596,2.1817,-0.6050) (5.8469,2.1817,-0.5843) (5.9341,2.1817,-0.5651) (6.0214,2.1817,-0.5473) (6.1087,2.1817,-0.5308) (6.1959,2.1817,-0.5152) (6.2832,2.1817,-0.5000)
(0.0000,2.2689,-0.5000) (0.0873,2.2689,-0.5129) (0.1745,2.2689,-0.5244) (0.2618,2.2689,-0.5333) (0.3491,2.2689,-0.5386) (0.4363,2.2689,-0.5393) (0.5236,2.2689,-0.5350) (0.6109,2.2689,-0.5256) (0.6981,2.2689,-0.5113) (0.7854,2.2689,-0.4927) (0.8727,2.2689,-0.4707) (0.9599,2.2689,-0.4465) (1.0472,2.2689,-0.4215) (1.1345,2.2689,-0.3970) (1.2217,2.2689,-0.3746) (1.3090,2.2689,-0.3556) (1.3963,2.2689,-0.3411) (1.4835,2.2689,-0.3321) (1.5708,2.2689,-0.3290) (1.6581,2.2689,-0.3321) (1.7453,2.2689,-0.3411) (1.8326,2.2689,-0.3556) (1.9199,2.2689,-0.3746) (2.0071,2.2689,-0.3970) (2.0944,2.2689,-0.4215) (2.1817,2.2689,-0.4465) (2.2689,2.2689,-0.4707) (2.3562,2.2689,-0.4927) (2.4435,2.2689,-0.5113) (2.5307,2.2689,-0.5256) (2.6180,2.2689,-0.5350) (2.7053,2.2689,-0.5393) (2.7925,2.2689,-0.5386) (2.8798,2.2689,-0.5333) (2.9671,2.2689,-0.5244) (3.0543,2.2689,-0.5129) (3.1416,2.2689,-0.5000) (3.2289,2.2689,-0.4871) (3.3161,2.2689,-0.4756) (3.4034,2.2689,-0.4667) (3.4907,2.2689,-0.4614) (3.5779,2.2689,-0.4607) (3.6652,2.2689,-0.4650) (3.7525,2.2689,-0.4744) (3.8397,2.2689,-0.4887) (3.9270,2.2689,-0.5073) (4.0143,2.2689,-0.5293) (4.1015,2.2689,-0.5535) (4.1888,2.2689,-0.5785) (4.2761,2.2689,-0.6030) (4.3633,2.2689,-0.6254) (4.4506,2.2689,-0.6444) (4.5379,2.2689,-0.6589) (4.6251,2.2689,-0.6679) (4.7124,2.2689,-0.6710) (4.7997,2.2689,-0.6679) (4.8869,2.2689,-0.6589) (4.9742,2.2689,-0.6444) (5.0615,2.2689,-0.6254) (5.1487,2.2689,-0.6030) (5.2360,2.2689,-0.5785) (5.3233,2.2689,-0.5535) (5.4105,2.2689,-0.5293) (5.4978,2.2689,-0.5073) (5.5851,2.2689,-0.4887) (5.6723,2.2689,-0.4744) (5.7596,2.2689,-0.4650) (5.8469,2.2689,-0.4607) (5.9341,2.2689,-0.4614) (6.0214,2.2689,-0.4667) (6.1087,2.2689,-0.4756) (6.1959,2.2689,-0.4871) (6.2832,2.2689,-0.5000)
(0.0000,2.3562,-0.5000) (0.0873,2.3562,-0.5432) (0.1745,2.3562,-0.5842) (0.2618,2.3562,-0.6207) (0.3491,2.3562,-0.6510) (0.4363,2.3562,-0.6736) (0.5236,2.3562,-0.6875) (0.6109,2.3562,-0.6924) (0.6981,2.3562,-0.6886) (0.7854,2.3562,-0.6768) (0.8727,2.3562,-0.6583) (0.9599,2.3562,-0.6347) (1.0472,2.3562,-0.6083) (1.1345,2.3562,-0.5809) (1.2217,2.3562,-0.5550) (1.3090,2.3562,-0.5324) (1.3963,2.3562,-0.5148) (1.4835,2.3562,-0.5038) (1.5708,2.3562,-0.5000) (1.6581,2.3562,-0.5038) (1.7453,2.3562,-0.5148) (1.8326,2.3562,-0.5324) (1.9199,2.3562,-0.5550) (2.0071,2.3562,-0.5809) (2.0944,2.3562,-0.6083) (2.1817,2.3562,-0.6347) (2.2689,2.3562,-0.6583) (2.3562,2.3562,-0.6768) (2.4435,2.3562,-0.6886) (2.5307,2.3562,-0.6924) (2.6180,2.3562,-0.6875) (2.7053,2.3562,-0.6736) (2.7925,2.3562,-0.6510) (2.8798,2.3562,-0.6207) (2.9671,2.3562,-0.5842) (3.0543,2.3562,-0.5432) (3.1416,2.3562,-0.5000) (3.2289,2.3562,-0.4568) (3.3161,2.3562,-0.4158) (3.4034,2.3562,-0.3793) (3.4907,2.3562,-0.3490) (3.5779,2.3562,-0.3264) (3.6652,2.3562,-0.3125) (3.7525,2.3562,-0.3076) (3.8397,2.3562,-0.3114) (3.9270,2.3562,-0.3232) (4.0143,2.3562,-0.3417) (4.1015,2.3562,-0.3653) (4.1888,2.3562,-0.3917) (4.2761,2.3562,-0.4191) (4.3633,2.3562,-0.4450) (4.4506,2.3562,-0.4676) (4.5379,2.3562,-0.4852) (4.6251,2.3562,-0.4962) (4.7124,2.3562,-0.5000) (4.7997,2.3562,-0.4962) (4.8869,2.3562,-0.4852) (4.9742,2.3562,-0.4676) (5.0615,2.3562,-0.4450) (5.1487,2.3562,-0.4191) (5.2360,2.3562,-0.3917) (5.3233,2.3562,-0.3653) (5.4105,2.3562,-0.3417) (5.4978,2.3562,-0.3232) (5.5851,2.3562,-0.3114) (5.6723,2.3562,-0.3076) (5.7596,2.3562,-0.3125) (5.8469,2.3562,-0.3264) (5.9341,2.3562,-0.3490) (6.0214,2.3562,-0.3793) (6.1087,2.3562,-0.4158) (6.1959,2.3562,-0.4568) (6.2832,2.3562,-0.5000)
(0.0000,2.4435,-0.5000) (0.0873,2.4435,-0.5723) (0.1745,2.4435,-0.6414) (0.2618,2.4435,-0.7045) (0.3491,2.4435,-0.7588) (0.4363,2.4435,-0.8026) (0.5236,2.4435,-0.8343) (0.6109,2.4435,-0.8534) (0.6981,2.4435,-0.8602) (0.7854,2.4435,-0.8555) (0.8727,2.4435,-0.8410) (0.9599,2.4435,-0.8189) (1.0472,2.4435,-0.7917) (1.1345,2.4435,-0.7624) (1.2217,2.4435,-0.7336) (1.3090,2.4435,-0.7081) (1.3963,2.4435,-0.6881) (1.4835,2.4435,-0.6754) (1.5708,2.4435,-0.6710) (1.6581,2.4435,-0.6754) (1.7453,2.4435,-0.6881) (1.8326,2.4435,-0.7081) (1.9199,2.4435,-0.7336) (2.0071,2.4435,-0.7624) (2.0944,2.4435,-0.7917) (2.1817,2.4435,-0.8189) (2.2689,2.4435,-0.8410) (2.3562,2.4435,-0.8555) (2.4435,2.4435,-0.8602) (2.5307,2.4435,-0.8534) (2.6180,2.4435,-0.8343) (2.7053,2.4435,-0.8026) (2.7925,2.4435,-0.7588) (2.8798,2.4435,-0.7045) (2.9671,2.4435,-0.6414) (3.0543,2.4435,-0.5723) (3.1416,2.4435,-0.5000) (3.2289,2.4435,-0.4277) (3.3161,2.4435,-0.3586) (3.4034,2.4435,-0.2955) (3.4907,2.4435,-0.2412) (3.5779,2.4435,-0.1974) (3.6652,2.4435,-0.1657) (3.7525,2.4435,-0.1466) (3.8397,2.4435,-0.1398) (3.9270,2.4435,-0.1445) (4.0143,2.4435,-0.1590) (4.1015,2.4435,-0.1811) (4.1888,2.4435,-0.2083) (4.2761,2.4435,-0.2376) (4.3633,2.4435,-0.2664) (4.4506,2.4435,-0.2919) (4.5379,2.4435,-0.3119) (4.6251,2.4435,-0.3246) (4.7124,2.4435,-0.3290) (4.7997,2.4435,-0.3246) (4.8869,2.4435,-0.3119) (4.9742,2.4435,-0.2919) (5.0615,2.4435,-0.2664) (5.1487,2.4435,-0.2376) (5.2360,2.4435,-0.2083) (5.3233,2.4435,-0.1811) (5.4105,2.4435,-0.1590) (5.4978,2.4435,-0.1445) (5.5851,2.4435,-0.1398) (5.6723,2.4435,-0.1466) (5.7596,2.4435,-0.1657) (5.8469,2.4435,-0.1974) (5.9341,2.4435,-0.2412) (6.0214,2.4435,-0.2955) (6.1087,2.4435,-0.3586) (6.1959,2.4435,-0.4277) (6.2832,2.4435,-0.5000)
(0.0000,2.5307,-0.5000) (0.0873,2.5307,-0.5964) (0.1745,2.5307,-0.6891) (0.2618,2.5307,-0.7743) (0.3491,2.5307,-0.8489) (0.4363,2.5307,-0.9105) (0.5236,2.5307,-0.9574) (0.6109,2.5307,-0.9889) (0.6981,2.5307,-1.0050) (0.7854,2.5307,-1.0070) (0.8727,2.5307,-0.9966) (0.9599,2.5307,-0.9765) (1.0472,2.5307,-0.9496) (1.1345,2.5307,-0.9194) (1.2217,2.5307,-0.8890) (1.3090,2.5307,-0.8616) (1.3963,2.5307,-0.8400) (1.4835,2.5307,-0.8262) (1.5708,2.5307,-0.8214) (1.6581,2.5307,-0.8262) (1.7453,2.5307,-0.8400) (1.8326,2.5307,-0.8616) (1.9199,2.5307,-0.8890) (2.0071,2.5307,-0.9194) (2.0944,2.5307,-0.9496) (2.1817,2.5307,-0.9765) (2.2689,2.5307,-0.9966) (2.3562,2.5307,-1.0070) (2.4435,2.5307,-1.0050) (2.5307,2.5307,-0.9889) (2.6180,2.5307,-0.9574) (2.7053,2.5307,-0.9105) (2.7925,2.5307,-0.8489) (2.8798,2.5307,-0.7743) (2.9671,2.5307,-0.6891) (3.0543,2.5307,-0.5964) (3.1416,2.5307,-0.5000) (3.2289,2.5307,-0.4036) (3.3161,2.5307,-0.3109) (3.4034,2.5307,-0.2257) (3.4907,2.5307,-0.1511) (3.5779,2.5307,-0.0895) (3.6652,2.5307,-0.0426) (3.7525,2.5307,-0.0111) (3.8397,2.5307,0.0050) (3.9270,2.5307,0.0070) (4.0143,2.5307,-0.0034) (4.1015,2.5307,-0.0235) (4.1888,2.5307,-0.0504) (4.2761,2.5307,-0.0806) (4.3633,2.5307,-0.1110) (4.4506,2.5307,-0.1384) (4.5379,2.5307,-0.1600) (4.6251,2.5307,-0.1738) (4.7124,2.5307,-0.1786) (4.7997,2.5307,-0.1738) (4.8869,2.5307,-0.1600) (4.9742,2.5307,-0.1384) (5.0615,2.5307,-0.1110) (5.1487,2.5307,-0.0806) (5.2360,2.5307,-0.0504) (5.3233,2.5307,-0.0235) (5.4105,2.5307,-0.0034) (5.4978,2.5307,0.0070) (5.5851,2.5307,0.0050) (5.6723,2.5307,-0.0111) (5.7596,2.5307,-0.0426) (5.8469,2.5307,-0.0895) (5.9341,2.5307,-0.1511) (6.0214,2.5307,-0.2257) (6.1087,2.5307,-0.3109) (6.1959,2.5307,-0.4036) (6.2832,2.5307,-0.5000)
(0.0000,2.6180,-0.5000) (0.0873,2.6180,-0.6126) (0.1745,2.6180,-0.7210) (0.2618,2.6180,-0.8212) (0.3491,2.6180,-0.9096) (0.4363,2.6180,-0.9836) (0.5236,2.6180,-1.0413) (0.6109,2.6180,-1.0817) (0.6981,2.6180,-1.1050) (0.7854,2.6180,-1.1124) (0.8727,2.6180,-1.1058) (0.9599,2.6180,-1.0881) (1.0472,2.6180,-1.0625) (1.1345,2.6180,-1.0326) (1.2217,2.6180,-1.0021) (1.3090,2.6180,-0.9743) (1.3963,2.6180,-0.9522) (1.4835,2.6180,-0.9379) (1.5708,2.6180,-0.9330) (1.6581,2.6180,-0.9379) (1.7453,2.6180,-0.9522) (1.8326,2.6180,-0.9743) (1.9199,2.6180,-1.0021) (2.0071,2.6180,-1.0326) (2.0944,2.6180,-1.0625) (2.1817,2.6180,-1.0881) (2.2689,2.6180,-1.1058) (2.3562,2.6180,-1.1124) (2.4435,2.6180,-1.1050) (2.5307,2.6180,-1.0817) (2.6180,2.6180,-1.0413) (2.7053,2.6180,-0.9836) (2.7925,2.6180,-0.9096) (2.8798,2.6180,-0.8212) (2.9671,2.6180,-0.7210) (3.0543,2.6180,-0.6126) (3.1416,2.6180,-0.5000) (3.2289,2.6180,-0.3874) (3.3161,2.6180,-0.2790) (3.4034,2.6180,-0.1788) (3.4907,2.6180,-0.0904) (3.5779,2.6180,-0.0164) (3.6652,2.6180,0.0413) (3.7525,2.6180,0.0817) (3.8397,2.6180,0.1050) (3.9270,2.6180,0.1124) (4.0143,2.6180,0.1058) (4.1015,2.6180,0.0881) (4.1888,2.6180,0.0625) (4.2761,2.6180,0.0326) (4.3633,2.6180,0.0021) (4.4506,2.6180,-0.0257) (4.5379,2.6180,-0.0478) (4.6251,2.6180,-0.0621) (4.7124,2.6180,-0.0670) (4.7997,2.6180,-0.0621) (4.8869,2.6180,-0.0478) (4.9742,2.6180,-0.0257) (5.0615,2.6180,0.0021) (5.1487,2.6180,0.0326) (5.2360,2.6180,0.0625) (5.3233,2.6180,0.0881) (5.4105,2.6180,0.1058) (5.4978,2.6180,0.1124) (5.5851,2.6180,0.1050) (5.6723,2.6180,0.0817) (5.7596,2.6180,0.0413) (5.8469,2.6180,-0.0164) (5.9341,2.6180,-0.0904) (6.0214,2.6180,-0.1788) (6.1087,2.6180,-0.2790) (6.1959,2.6180,-0.3874) (6.2832,2.6180,-0.5000)
(0.0000,2.7053,-0.5000) (0.0873,2.7053,-0.6186) (0.1745,2.7053,-0.7329) (0.2618,2.7053,-0.8388) (0.3491,2.7053,-0.9328) (0.4363,2.7053,-1.0120) (0.5236,2.7053,-1.0745) (0.6109,2.7053,-1.1194) (0.6981,2.7053,-1.1467) (0.7854,2.7053,-1.1577) (0.8727,2.7053,-1.1543) (0.9599,2.7053,-1.1393) (1.0472,2.7053,-1.1160) (1.1345,2.7053,-1.0880) (1.2217,2.7053,-1.0589) (1.3090,2.7053,-1.0323) (1.3963,2.7053,-1.0109) (1.4835,2.7053,-0.9972) (1.5708,2.7053,-0.9924) (1.6581,2.7053,-0.9972) (1.7453,2.7053,-1.0109) (1.8326,2.7053,-1.0323) (1.9199,2.7053,-1.0589) (2.0071,2.7053,-1.0880) (2.0944,2.7053,-1.1160) (2.1817,2.7053,-1.1393) (2.2689,2.7053,-1.1543) (2.3562,2.7053,-1.1577) (2.4435,2.7053,-1.1467) (2.5307,2.7053,-1.1194) (2.6180,2.7053,-1.0745) (2.7053,2.7053,-1.0120) (2.7925,2.7053,-0.9328) (2.8798,2.7053,-0.8388) (2.9671,2.7053,-0.7329) (3.0543,2.7053,-0.6186) (3.1416,2.7053,-0.5000) (3.2289,2.7053,-0.3814) (3.3161,2.7053,-0.2671) (3.4034,2.7053,-0.1612) (3.4907,2.7053,-0.0672) (3.5779,2.7053,0.0120) (3.6652,2.7053,0.0745) (3.7525,2.7053,0.1194) (3.8397,2.7053,0.1467) (3.9270,2.7053,0.1577) (4.0143,2.7053,0.1543) (4.1015,2.7053,0.1393) (4.1888,2.7053,0.1160) (4.2761,2.7053,0.0880) (4.3633,2.7053,0.0589) (4.4506,2.7053,0.0323) (4.5379,2.7053,0.0109) (4.6251,2.7053,-0.0028) (4.7124,2.7053,-0.0076) (4.7997,2.7053,-0.0028) (4.8869,2.7053,0.0109) (4.9742,2.7053,0.0323) (5.0615,2.7053,0.0589) (5.1487,2.7053,0.0880) (5.2360,2.7053,0.1160) (5.3233,2.7053,0.1393) (5.4105,2.7053,0.1543) (5.4978,2.7053,0.1577) (5.5851,2.7053,0.1467) (5.6723,2.7053,0.1194) (5.7596,2.7053,0.0745) (5.8469,2.7053,0.0120) (5.9341,2.7053,-0.0672) (6.0214,2.7053,-0.1612) (6.1087,2.7053,-0.2671) (6.1959,2.7053,-0.3814) (6.2832,2.7053,-0.5000)
(0.0000,2.7925,-0.5000) (0.0873,2.7925,-0.6133) (0.1745,2.7925,-0.7226) (0.2618,2.7925,-0.8240) (0.3491,2.7925,-0.9142) (0.4363,2.7925,-0.9906) (0.5236,2.7925,-1.0514) (0.6109,2.7925,-1.0956) (0.6981,2.7925,-1.1235) (0.7854,2.7925,-1.1359) (0.8727,2.7925,-1.1348) (0.9599,2.7925,-1.1227) (1.0472,2.7925,-1.1026) (1.1345,2.7925,-1.0780) (1.2217,2.7925,-1.0522) (1.3090,2.7925,-1.0283) (1.3963,2.7925,-1.0091) (1.4835,2.7925,-0.9967) (1.5708,2.7925,-0.9924) (1.6581,2.7925,-0.9967) (1.7453,2.7925,-1.0091) (1.8326,2.7925,-1.0283) (1.9199,2.7925,-1.0522) (2.0071,2.7925,-1.0780) (2.0944,2.7925,-1.1026) (2.1817,2.7925,-1.1227) (2.2689,2.7925,-1.1348) (2.3562,2.7925,-1.1359) (2.4435,2.7925,-1.1235) (2.5307,2.7925,-1.0956) (2.6180,2.7925,-1.0514) (2.7053,2.7925,-0.9906) (2.7925,2.7925,-0.9142) (2.8798,2.7925,-0.8240) (2.9671,2.7925,-0.7226) (3.0543,2.7925,-0.6133) (3.1416,2.7925,-0.5000) (3.2289,2.7925,-0.3867) (3.3161,2.7925,-0.2774) (3.4034,2.7925,-0.1760) (3.4907,2.7925,-0.0858) (3.5779,2.7925,-0.0094) (3.6652,2.7925,0.0514) (3.7525,2.7925,0.0956) (3.8397,2.7925,0.1235) (3.9270,2.7925,0.1359) (4.0143,2.7925,0.1348) (4.1015,2.7925,0.1227) (4.1888,2.7925,0.1026) (4.2761,2.7925,0.0780) (4.3633,2.7925,0.0522) (4.4506,2.7925,0.0283) (4.5379,2.7925,0.0091) (4.6251,2.7925,-0.0033) (4.7124,2.7925,-0.0076) (4.7997,2.7925,-0.0033) (4.8869,2.7925,0.0091) (4.9742,2.7925,0.0283) (5.0615,2.7925,0.0522) (5.1487,2.7925,0.0780) (5.2360,2.7925,0.1026) (5.3233,2.7925,0.1227) (5.4105,2.7925,0.1348) (5.4978,2.7925,0.1359) (5.5851,2.7925,0.1235) (5.6723,2.7925,0.0956) (5.7596,2.7925,0.0514) (5.8469,2.7925,-0.0094) (5.9341,2.7925,-0.0858) (6.0214,2.7925,-0.1760) (6.1087,2.7925,-0.2774) (6.1959,2.7925,-0.3867) (6.2832,2.7925,-0.5000)
(0.0000,2.8798,-0.5000) (0.0873,2.8798,-0.5968) (0.1745,2.8798,-0.6902) (0.2618,2.8798,-0.7770) (0.3491,2.8798,-0.8544) (0.4363,2.8798,-0.9201) (0.5236,2.8798,-0.9726) (0.6109,2.8798,-1.0112) (0.6981,2.8798,-1.0360) (0.7854,2.8798,-1.0477) (0.8727,2.8798,-1.0479) (0.9599,2.8798,-1.0388) (1.0472,2.8798,-1.0229) (1.1345,2.8798,-1.0030) (1.2217,2.8798,-0.9820) (1.3090,2.8798,-0.9625) (1.3963,2.8798,-0.9467) (1.4835,2.8798,-0.9365) (1.5708,2.8798,-0.9330) (1.6581,2.8798,-0.9365) (1.7453,2.8798,-0.9467) (1.8326,2.8798,-0.9625) (1.9199,2.8798,-0.9820) (2.0071,2.8798,-1.0030) (2.0944,2.8798,-1.0229) (2.1817,2.8798,-1.0388) (2.2689,2.8798,-1.0479) (2.3562,2.8798,-1.0477) (2.4435,2.8798,-1.0360) (2.5307,2.8798,-1.0112) (2.6180,2.8798,-0.9726) (2.7053,2.8798,-0.9201) (2.7925,2.8798,-0.8544) (2.8798,2.8798,-0.7770) (2.9671,2.8798,-0.6902) (3.0543,2.8798,-0.5968) (3.1416,2.8798,-0.5000) (3.2289,2.8798,-0.4032) (3.3161,2.8798,-0.3098) (3.4034,2.8798,-0.2230) (3.4907,2.8798,-0.1456) (3.5779,2.8798,-0.0799) (3.6652,2.8798,-0.0274) (3.7525,2.8798,0.0112) (3.8397,2.8798,0.0360) (3.9270,2.8798,0.0477) (4.0143,2.8798,0.0479) (4.1015,2.8798,0.0388) (4.1888,2.8798,0.0229) (4.2761,2.8798,0.0030) (4.3633,2.8798,-0.0180) (4.4506,2.8798,-0.0375) (4.5379,2.8798,-0.0533) (4.6251,2.8798,-0.0635) (4.7124,2.8798,-0.0670) (4.7997,2.8798,-0.0635) (4.8869,2.8798,-0.0533) (4.9742,2.8798,-0.0375) (5.0615,2.8798,-0.0180) (5.1487,2.8798,0.0030) (5.2360,2.8798,0.0229) (5.3233,2.8798,0.0388) (5.4105,2.8798,0.0479) (5.4978,2.8798,0.0477) (5.5851,2.8798,0.0360) (5.6723,2.8798,0.0112) (5.7596,2.8798,-0.0274) (5.8469,2.8798,-0.0799) (5.9341,2.8798,-0.1456) (6.0214,2.8798,-0.2230) (6.1087,2.8798,-0.3098) (6.1959,2.8798,-0.4032) (6.2832,2.8798,-0.5000)
(0.0000,2.9671,-0.5000) (0.0873,2.9671,-0.5706) (0.1745,2.9671,-0.6387) (0.2618,2.9671,-0.7021) (0.3491,2.9671,-0.7586) (0.4363,2.9671,-0.8068) (0.5236,2.9671,-0.8453) (0.6109,2.9671,-0.8739) (0.6981,2.9671,-0.8923) (0.7854,2.9671,-0.9014) (0.8727,2.9671,-0.9021) (0.9599,2.9671,-0.8960) (1.0472,2.9671,-0.8849) (1.1345,2.9671,-0.8710) (1.2217,2.9671,-0.8561) (1.3090,2.9671,-0.8423) (1.3963,2.9671,-0.8311) (1.4835,2.9671,-0.8239) (1.5708,2.9671,-0.8214) (1.6581,2.9671,-0.8239) (1.7453,2.9671,-0.8311) (1.8326,2.9671,-0.8423) (1.9199,2.9671,-0.8561) (2.0071,2.9671,-0.8710) (2.0944,2.9671,-0.8849) (2.1817,2.9671,-0.8960) (2.2689,2.9671,-0.9021) (2.3562,2.9671,-0.9014) (2.4435,2.9671,-0.8923) (2.5307,2.9671,-0.8739) (2.6180,2.9671,-0.8453) (2.7053,2.9671,-0.8068) (2.7925,2.9671,-0.7586) (2.8798,2.9671,-0.7021) (2.9671,2.9671,-0.6387) (3.0543,2.9671,-0.5706) (3.1416,2.9671,-0.5000) (3.2289,2.9671,-0.4294) (3.3161,2.9671,-0.3613) (3.4034,2.9671,-0.2979) (3.4907,2.9671,-0.2414) (3.5779,2.9671,-0.1932) (3.6652,2.9671,-0.1547) (3.7525,2.9671,-0.1261) (3.8397,2.9671,-0.1077) (3.9270,2.9671,-0.0986) (4.0143,2.9671,-0.0979) (4.1015,2.9671,-0.1040) (4.1888,2.9671,-0.1151) (4.2761,2.9671,-0.1290) (4.3633,2.9671,-0.1439) (4.4506,2.9671,-0.1577) (4.5379,2.9671,-0.1689) (4.6251,2.9671,-0.1761) (4.7124,2.9671,-0.1786) (4.7997,2.9671,-0.1761) (4.8869,2.9671,-0.1689) (4.9742,2.9671,-0.1577) (5.0615,2.9671,-0.1439) (5.1487,2.9671,-0.1290) (5.2360,2.9671,-0.1151) (5.3233,2.9671,-0.1040) (5.4105,2.9671,-0.0979) (5.4978,2.9671,-0.0986) (5.5851,2.9671,-0.1077) (5.6723,2.9671,-0.1261) (5.7596,2.9671,-0.1547) (5.8469,2.9671,-0.1932) (5.9341,2.9671,-0.2414) (6.0214,2.9671,-0.2979) (6.1087,2.9671,-0.3613) (6.1959,2.9671,-0.4294) (6.2832,2.9671,-0.5000)
(0.0000,3.0543,-0.5000) (0.0873,3.0543,-0.5372) (0.1745,3.0543,-0.5731) (0.2618,3.0543,-0.6065) (0.3491,3.0543,-0.6364) (0.4363,3.0543,-0.6618) (0.5236,3.0543,-0.6822) (0.6109,3.0543,-0.6973) (0.6981,3.0543,-0.7072) (0.7854,3.0543,-0.7121) (0.8727,3.0543,-0.7126) (0.9599,3.0543,-0.7096) (1.0472,3.0543,-0.7039) (1.1345,3.0543,-0.6967) (1.2217,3.0543,-0.6890) (1.3090,3.0543,-0.6819) (1.3963,3.0543,-0.6761) (1.4835,3.0543,-0.6723) (1.5708,3.0543,-0.6710) (1.6581,3.0543,-0.6723) (1.7453,3.0543,-0.6761) (1.8326,3.0543,-0.6819) (1.9199,3.0543,-0.6890) (2.0071,3.0543,-0.6967) (2.0944,3.0543,-0.7039) (2.1817,3.0543,-0.7096) (2.2689,3.0543,-0.7126) (2.3562,3.0543,-0.7121) (2.4435,3.0543,-0.7072) (2.5307,3.0543,-0.6973) (2.6180,3.0543,-0.6822) (2.7053,3.0543,-0.6618) (2.7925,3.0543,-0.6364) (2.8798,3.0543,-0.6065) (2.9671,3.0543,-0.5731) (3.0543,3.0543,-0.5372) (3.1416,3.0543,-0.5000) (3.2289,3.0543,-0.4628) (3.3161,3.0543,-0.4269) (3.4034,3.0543,-0.3935) (3.4907,3.0543,-0.3636) (3.5779,3.0543,-0.3382) (3.6652,3.0543,-0.3178) (3.7525,3.0543,-0.3027) (3.8397,3.0543,-0.2928) (3.9270,3.0543,-0.2879) (4.0143,3.0543,-0.2874) (4.1015,3.0543,-0.2904) (4.1888,3.0543,-0.2961) (4.2761,3.0543,-0.3033) (4.3633,3.0543,-0.3110) (4.4506,3.0543,-0.3181) (4.5379,3.0543,-0.3239) (4.6251,3.0543,-0.3277) (4.7124,3.0543,-0.3290) (4.7997,3.0543,-0.3277) (4.8869,3.0543,-0.3239) (4.9742,3.0543,-0.3181) (5.0615,3.0543,-0.3110) (5.1487,3.0543,-0.3033) (5.2360,3.0543,-0.2961) (5.3233,3.0543,-0.2904) (5.4105,3.0543,-0.2874) (5.4978,3.0543,-0.2879) (5.5851,3.0543,-0.2928) (5.6723,3.0543,-0.3027) (5.7596,3.0543,-0.3178) (5.8469,3.0543,-0.3382) (5.9341,3.0543,-0.3636) (6.0214,3.0543,-0.3935) (6.1087,3.0543,-0.4269) (6.1959,3.0543,-0.4628) (6.2832,3.0543,-0.5000)
(0.0000,3.1416,-0.5000) (0.0873,3.1416,-0.5000) (0.1745,3.1416,-0.5000) (0.2618,3.1416,-0.5000) (0.3491,3.1416,-0.5000) (0.4363,3.1416,-0.5000) (0.5236,3.1416,-0.5000) (0.6109,3.1416,-0.5000) (0.6981,3.1416,-0.5000) (0.7854,3.1416,-0.5000) (0.8727,3.1416,-0.5000) (0.9599,3.1416,-0.5000) (1.0472,3.1416,-0.5000) (1.1345,3.1416,-0.5000) (1.2217,3.1416,-0.5000) (1.3090,3.1416,-0.5000) (1.3963,3.1416,-0.5000) (1.4835,3.1416,-0.5000) (1.5708,3.1416,-0.5000) (1.6581,3.1416,-0.5000) (1.7453,3.1416,-0.5000) (1.8326,3.1416,-0.5000) (1.9199,3.1416,-0.5000) (2.0071,3.1416,-0.5000) (2.0944,3.1416,-0.5000) (2.1817,3.1416,-0.5000) (2.2689,3.1416,-0.5000) (2.3562,3.1416,-0.5000) (2.4435,3.1416,-0.5000) (2.5307,3.1416,-0.5000) (2.6180,3.1416,-0.5000) (2.7053,3.1416,-0.5000) (2.7925,3.1416,-0.5000) (2.8798,3.1416,-0.5000) (2.9671,3.1416,-0.5000) (3.0543,3.1416,-0.5000) (3.1416,3.1416,-0.5000) (3.2289,3.1416,-0.5000) (3.3161,3.1416,-0.5000) (3.4034,3.1416,-0.5000) (3.4907,3.1416,-0.5000) (3.5779,3.1416,-0.5000) (3.6652,3.1416,-0.5000) (3.7525,3.1416,-0.5000) (3.8397,3.1416,-0.5000) (3.9270,3.1416,-0.5000) (4.0143,3.1416,-0.5000) (4.1015,3.1416,-0.5000) (4.1888,3.1416,-0.5000) (4.2761,3.1416,-0.5000) (4.3633,3.1416,-0.5000) (4.4506,3.1416,-0.5000) (4.5379,3.1416,-0.5000) (4.6251,3.1416,-0.5000) (4.7124,3.1416,-0.5000) (4.7997,3.1416,-0.5000) (4.8869,3.1416,-0.5000) (4.9742,3.1416,-0.5000) (5.0615,3.1416,-0.5000) (5.1487,3.1416,-0.5000) (5.2360,3.1416,-0.5000) (5.3233,3.1416,-0.5000) (5.4105,3.1416,-0.5000) (5.4978,3.1416,-0.5000) (5.5851,3.1416,-0.5000) (5.6723,3.1416,-0.5000) (5.7596,3.1416,-0.5000) (5.8469,3.1416,-0.5000) (5.9341,3.1416,-0.5000) (6.0214,3.1416,-0.5000) (6.1087,3.1416,-0.5000) (6.1959,3.1416,-0.5000) (6.2832,3.1416,-0.5000)
};
\addplot3[only marks,mark=x,mark size=4pt,corange,very thick] coordinates {(0.7854,2.7053,0)};
\end{axis}
\end{tikzpicture}

\caption{Paysage $F_1(\gamma,\beta)$ du QUBO jouet (calcul exact sur une grille $73\times37$). La croix marque le minimum, $F_1\approx-1{,}16$, à comparer à l'énergie moyenne $-0{,}5$ de l'état initial et au fondamental $-2$. Le paysage est périodique ($2\pi$ en $\gamma$ car les énergies sont entières, $\pi$ en $\beta$) et non convexe : l'optimiseur classique a besoin de plusieurs points de départ.}\label{fig:paysage-qaoa}
\end{figure}

\begin{figure}[htbp]
\centering
% figure générée par tools/make_figures.py : jouet_proba_p1
\begin{tikzpicture}
\begin{axis}[width=0.8\linewidth,height=5cm,ybar,bar width=12pt,axis lines=left,ymin=0,ymax=0.4,
  xtick={0,...,7},xticklabels={$|000\rangle$,$|100\rangle$,$|010\rangle$,$|110\rangle$,$|001\rangle$,$|101\rangle$,$|011\rangle$,$|111\rangle$},
  ylabel={probabilité},xticklabel style={font=\small},
  nodes near coords={},enlarge x limits=0.08]
\addplot[fill=cblue,draw=none] coordinates {(0,0.04482) (1,0.1365) (2,0.2777) (3,0.04093) (4,0.1365) (5,0.3104) (6,0.04093) (7,0.01218)};
\addplot[corange,thick,sharp plot,mark=none,update limits=false] coordinates {(-0.5,0.125) (7.5,0.125)};
\end{axis}
\end{tikzpicture}

\caption{Distribution de mesure après QAOA $p=1$ aux angles optimaux. La solution $\ket{101}$ sort avec probabilité $0{,}31$, contre $1/8=0{,}125$ pour la superposition uniforme (trait orange). À $p=2$ et $p=3$, elle monte à $0{,}54$ puis $0{,}67$.}\label{fig:jouet-proba}
\end{figure}

\section{Résultats sur notre QUBO aéro (9 qubits)}
Le QUBO de traînée du chapitre~\ref{ch:qubo} (512 chaînes) a été traité par QAOA simulé exactement, avec une optimisation COBYLA. Pour $p=1$, on part de 8 points de départ aléatoires ; pour $p\ge2$, on part de l'interpolation des angles optimaux de $p-1$ (stratégie INTERP).

\begin{center}
\begin{tabular}{cccccc}
\toprule
$p$ & CNOT & $P(\text{optimum})$ & $P(\text{5 meilleurs})$ & ratio $r$ & shots moyens pour voir l'optimum\\
\midrule
1 & 72 & 0,018 & 0,088 & 0,867 & 56\\
2 & 144 & 0,062 & 0,218 & 0,933 & 16\\
3 & 216 & 0,097 & 0,302 & 0,952 & 10\\
4 & 288 & 0,127 & 0,384 & 0,966 & 8\\
5 & 360 & 0,142 & 0,436 & 0,972 & 7\\
\bottomrule
\end{tabular}
\end{center}
Le ratio d'approximation est $r=1-\expval{\tilde E}$, où $\tilde E=(E-E_{\min})/(E_{\max}-E_{\min})\in[0,1]$ ; il vaut 1 si l'on mesure toujours l'optimum. Le hasard uniforme donne $P(\text{optimum})=1/512\approx0{,}002$.

\begin{figure}[htbp]
\centering
% figure générée par tools/make_figures.py : qaoa_proba_rang
\begin{tikzpicture}
\begin{axis}[width=0.85\linewidth,height=5.8cm,axis lines=left,xmin=0.5,xmax=25.5,ymin=0,ymax=0.155,
  xlabel={rang de l'énergie QUBO (1 = minimum)},ylabel={probabilité de mesure},
  legend style={draw=none,font=\small,at={(0.98,0.98)},anchor=north east}]
\addplot[cblue,thick,mark=*,mark size=1.6pt] coordinates {(1,0.01771) (2,0.02594) (3,0.01501) (4,0.01886) (5,0.01022) (6,0.011) (7,0.03356) (8,0.0103) (9,0.007724) (10,0.01746) (11,0.007663) (12,0.01215) (13,0.01348) (14,0.005126) (15,0.009068) (16,0.004641) (17,0.02513) (18,0.007766) (19,0.007132) (20,0.00968) (21,0.01745) (22,0.004169) (23,0.005312) (24,0.005923) (25,0.0109)};\addlegendentry{$p=1$}
\addplot[corange,thick,mark=*,mark size=1.6pt] coordinates {(1,0.09662) (2,0.06801) (3,0.05075) (4,0.02962) (5,0.05687) (6,0.1012) (7,0.03086) (8,0.02036) (9,0.02548) (10,0.04532) (11,0.04281) (12,0.03432) (13,0.01705) (14,0.01967) (15,0.02236) (16,0.007503) (17,0.03259) (18,0.01114) (19,0.01694) (20,0.007158) (21,0.008521) (22,0.007297) (23,0.01132) (24,0.02091) (25,0.01149)};\addlegendentry{$p=3$}
\addplot[cgreen,thick,mark=*,mark size=1.6pt] coordinates {(1,0.142) (2,0.1067) (3,0.06809) (4,0.04338) (5,0.07616) (6,0.128) (7,0.04926) (8,0.02372) (9,0.02495) (10,0.02896) (11,0.06437) (12,0.03374) (13,0.01058) (14,0.0238) (15,0.01275) (16,0.005626) (17,0.00962) (18,0.005358) (19,0.02331) (20,0.004033) (21,0.005568) (22,0.003729) (23,0.005821) (24,0.01116) (25,0.002006)};\addlegendentry{$p=5$}
\addplot[cmuted,dashed,thick] coordinates {(0.5,0.00195) (25.5,0.00195)};\addlegendentry{hasard : $1/512$}
\end{axis}
\end{tikzpicture}

\caption{Probabilité de mesurer chaque chaîne, classée par énergie QUBO croissante (25 premières). QAOA concentre la probabilité sur le bas du spectre, d'autant plus que $p$ est grand. Pour $p=1$, le rang 2 dépasse le rang 1 : les deux profils sont quasi dégénérés (§\ref{sec:notre-qubo}).}\label{fig:qaoa-rang}
\end{figure}

\begin{figure}[htbp]
\centering
\begin{minipage}{0.49\linewidth}\centering% figure générée par tools/make_figures.py : qaoa_convergence_p
\begin{tikzpicture}
\begin{axis}[width=0.7\linewidth,height=5.4cm,axis lines=left,xmin=0.7,xmax=5.3,ymin=0,ymax=1,xtick={1,...,5},
  xlabel={profondeur $p$},ylabel={probabilité / ratio},
  legend style={draw=none,font=\small,at={(0.02,0.62)},anchor=west}]
\addplot[cblue,thick,mark=*] coordinates {(1,0.01771) (2,0.06215) (3,0.09662) (4,0.1269) (5,0.142)};\addlegendentry{$P(\text{optimum})$}
\addplot[corange,thick,mark=square*] coordinates {(1,0.08774) (2,0.2185) (3,0.3019) (4,0.3836) (5,0.4363)};\addlegendentry{$P(\text{5 meilleurs})$}
\addplot[cgreen,thick,mark=triangle*] coordinates {(1,0.8673) (2,0.9325) (3,0.9515) (4,0.9658) (5,0.972)};\addlegendentry{ratio d'approximation $r$}
\end{axis}
\end{tikzpicture}
\end{minipage}\hfill
\begin{minipage}{0.49\linewidth}\centering% figure générée par tools/make_figures.py : qaoa_angles
\begin{tikzpicture}
\begin{axis}[width=0.7\linewidth,height=5.2cm,axis lines=left,xmin=0.7,xmax=5.3,xtick={1,...,5},
  xlabel={couche $k$},ylabel={angle (rad)},
  legend style={draw=none,font=\small,at={(0.02,0.98)},anchor=north west}]
\addplot[cblue,thick,mark=*] coordinates {(1,0.5558) (2,1.153) (3,1.706) (4,2.239) (5,2.563)};\addlegendentry{$\gamma_k$ (coût)}
\addplot[corange,thick,mark=square*] coordinates {(1,0.6973) (2,0.5187) (3,0.3131) (4,0.2204) (5,0.1343)};\addlegendentry{$|\beta_k|$ (mélangeur)}
\end{axis}
\end{tikzpicture}
\end{minipage}
\caption{À gauche : qualité en fonction de la profondeur. À droite : angles optimaux pour $p=5$. $\gamma_k$ croît et $|\beta_k|$ décroît le long du circuit : c'est le profil d'un \emph{recuit} ($s$ croissant), retrouvé par l'optimiseur sans qu'on le lui impose.}\label{fig:qaoa-p}
\end{figure}

\begin{piege}[ (échelle de l'énergie)]
$\gamma$ multiplie $E$ dans une phase. Avec $E$ en counts de traînée (de $57$ à $117$), une variation de $\gamma$ de $0{,}1$ fait tourner les phases de plusieurs radians : le paysage $F_1$ devient une « tôle ondulée » où l'optimiseur se perd. Sans normalisation, on obtient $P(\text{optimum})=0{,}0045$ à $p=1$ ; en centrant et réduisant $E$ (\texttt{normalize}), on obtient $0{,}018$. Il faut toujours normaliser le QUBO avant QAOA.
\end{piege}

\section{Limites à connaître}
\begin{itemize}
\item \textbf{Optimisation classique difficile} : le paysage est non convexe, avec de nombreux minima locaux. Il faut des stratégies d'initialisation (INTERP, transfert d'angles d'une instance à une autre).
\item \textbf{Estimation bruitée} : sur un vrai processeur, $F_p$ est estimé avec $N$ shots, donc avec un bruit statistique en $1/\sqrt N$, plus le bruit matériel.
\item \textbf{Profondeur} : $2|E|p$ CNOT. Avec un graphe complet, la profondeur explose avec $n$, ce qui est incompatible avec le matériel bruité actuel au-delà de quelques couches.
\item \textbf{Plateaux arides} : pour des circuits profonds et larges, le gradient de $F_p$ peut devenir exponentiellement petit.
\item \textbf{Pas d'avantage prouvé} : QAOA à $p$ fixe n'est pas connu pour battre les meilleurs algorithmes classiques sur des QUBO généraux.
\end{itemize}

\exercices
\begin{exercice}
Généraliser l'exemple à un qubit à $H_C=h\,Z$ avec $h>0$. Montrer que $F_1(\gamma,\beta)=h\,\sin(2h\gamma)\,\sin2\beta$. Où sont les minima ? Que se passe-t-il pour le paysage si $h$ est très grand ? (Lien avec le piège de l'échelle.)
\end{exercice}
\begin{exercice}
Pour notre QUBO aéro (9 qubits, 36 couplages), compter les portes d'un circuit QAOA à $p=3$ : $H$, $R_Z$, CNOT, $R_X$. Avec $P(\text{optimum})=0{,}097$, combien de shots faut-il pour observer l'optimum au moins une fois avec une probabilité de 99~\% ?
\end{exercice}
